\subsection{CRITERIA FOR RATIONALITY}

For every subfield L of K, let \(\operatorname{End}_{\mathrm{L}}(\mathrm{K})\) denote the endomorphism ring of K considered as a left vector space over L; if L contains \(K'\) , \(\operatorname{End}_{\mathrm{L}}(\mathrm{K})\) is a subring of \(\operatorname{End}_{\mathrm{K}'}(\mathrm{K})\) . For every subset M of \(\operatorname{End}_{\mathrm{K}'}(\mathrm{K})\) , there exists a largest subfield L of K containing \(K'\) and such that M is contained in \(\operatorname{End}_{\mathrm{L}}(\mathrm{K})\) , namely the set of \(\xi \in K\) such that \(\phi(\xi\eta) = \xi.\phi(\eta)\) for all \(\eta \in K\) and all \(\phi \in M\) (it is immediately verified that this set is a ring and, on the other hand, replacing \(\eta\) by \(\xi^{-1}\eta\) in the preceding relation, we obtain \(\phi(\xi^{-1}\eta) = \xi^{-1}.\phi(\eta)\) when \(\xi \neq 0\) ). We shall call this field the centralizer of M in K and denote it by \(\chi(\mathcal{M})\) . Now let V be a right vector K-space with a \(K'\) -structure \(V'\) . For all \(\phi \in \operatorname{End}_{\mathrm{K}'}(\mathrm{K})\) , there exists one and only one endomorphism \(\phi_{V}\) of the Z-module V such that \(\phi_{\mathbf{V}}(x'.\xi) = x'.\phi(\xi)\) for \(x' \in \mathbf{V}'\) and \(\xi \in \mathbf{K}\) : for, in no. 1, a Z-isomorphism \(\lambda\) of \(\mathbf{V}' \otimes_{\mathbb{K}'} \mathbf{K}\) onto V was defined which maps \(x' \otimes \xi\) to \(x'\) . \(\xi\) and \(\phi_{\mathbf{V}}\) is necessarily equal to \(\lambda \circ (1_{\mathbf{V}'} \otimes \phi) \circ \lambda^{-1}\) .

THEOREM 1. Let \(\mathcal{M}\) be a subset of \(\operatorname{End}_{\mathbb{K}'}(\mathbf{K})\) and \(\mathbf{L} = \chi(\mathcal{M})\) the subfield of \(\mathbf{K}\) the centralizer of \(\mathcal{M}\) .

\begin{enumerate}
\def\labelenumi{(\roman{enumi})}
\item
  Let V be a right vector K-space with a K'-structure. For a vector \(x \in V\) to be rational over L, it is necessary and sufficient that \(\phi_{\mathbf{v}}(x.\eta) = x.\phi(\eta)\) for all \(\phi \in M\) and all \(\eta \in K\) . For a vector sub-K-space W of V to be rational over L, it is necessary and sufficient that \(\phi_{\mathbf{v}}(\mathbf{W}) \subset \mathbf{W}\) for all \(\phi \in M\) .
\item
  Let \(V_{1}\) , \(V_{2}\) be two right vector K-spaces each with a \(K^{\prime}\) -structure. For a K-linear mapping f of \(V_{1}\) into \(V_{2}\) to be rational over L, it is necessary and sufficient that \(f(\phi_{\mathrm{V}_{1}}(x_{1})) = \phi_{\mathrm{V}_{2}}(f(x_{1}))\) for all \(x_{1} \in V_{1}\) and all \(\phi \in M\) .
\end{enumerate}

We first prove assertion (i) for \(x\) . Let \(B\) be a basis of \(V\) rational over \(K'\) and write \(x = \sum_{b \in B} b \cdot \xi_b\) ; then, for \(\phi \in \mathcal{M}\) and \(\eta \in K\) ,

\[
\phi_ {V} (x. \eta) - x. \phi (\eta) = \sum_ {b \in B} b. (\phi (\xi_ {b} \eta) - \xi_ {b}. \phi (\eta))
\]

and therefore, the relations

\[
" \text { for   all } \phi \in \mathcal {M} \text { and   all } \eta \in K, \phi_ {V} (x. \eta) = x. \phi (\eta)"
\]

and

\[
" \text { for   all } \phi \in \mathcal {M}, \text { all } b \in \mathrm{B} \text { and   all } \eta \in \mathrm{K}, \phi (\xi_ {b} \eta) = \xi_ {b}. \phi (\eta)"
\]

are equivalent. The second of these relations means that for all \(b \in \mathbf{B}\) , \(\xi_b \in \chi(\mathcal{M})\) , which proves the first assertion of (i).

We next prove (ii). For \(f\) to be rational over \(L\) , it is necessary and sufficient that, for all \(x_1' \in V_1\) rational over \(K', f(x_1')\) is a vector in \(V_2\) rational over \(L\) ; this will imply that \(f(x_1)\) is rational over \(L\) for every vector \(x_1\) of \(V_1\) rational over \(L\) , such a vector being a linear combination with coefficients in \(L\) of vectors rational over \(K'\) . The above condition is equivalent, by the first part of the argument, to the relation

\[
f (x _ {1} ^ {\prime}) \cdot \phi (\eta) = \phi_ {\mathrm{v} _ {2}} (f (x _ {1} ^ {\prime}) \cdot \eta) \quad \text { for } \phi \in \mathcal {M} \text { and } \eta \in \mathrm{K}\tag{4}
\]

which may also be written

\[
f (\phi_ {\mathrm{v} _ {1}} (x _ {1} ^ {\prime}. \eta)) = \phi_ {\mathrm{v} _ {2}} (f (x _ {1} ^ {\prime}. \eta)) \quad \text { for } \phi \in \mathcal {M} \text { and } \eta \in \mathrm{K}.\tag{5}
\]

As every element of \(V_{1}\) is a linear combination with coefficients in K of elements of \(V_{1}\) rational over \(K'\) , condition (5) is equivalent to

\[
f (\phi_ {\mathrm{v} _ {1}} (x _ {1})) = \phi_ {\mathrm{v} _ {2}} (f (x _ {1}))
\]

for all \(x_{1} \in V_{1}\) and all \(\phi \in \mathcal{M}\) .

Finally, to prove the second assertion in (i), we use no. 6, Lemma 1: W is the graph of a K-linear mapping \(g \colon \mathbb{W}_1 \to \mathbb{W}_2\) and W is rational over L if and only if the mapping \(g\) is rational over L (no. 3, Proposition 4). By (ii), for \(g\) to be rational over L, it is necessary and sufficient that \(g(\phi_{\mathbb{W}_1}(x_1)) = \phi_{\mathbb{W}_2}(g(x_1))\) for all \(x_1 \in \mathbb{W}_1\) and all \(\phi \in \mathcal{M}\) ; as \(\phi_{\mathbb{V}} = \phi_{\mathbb{W}_1} \times \phi_{\mathbb{W}_2}\) , the above condition means that the graph W of \(g\) is stable under \(\phi_{\mathbb{V}}\) for all \(\phi \in \mathcal{M}\) .

\section{AFFINE SPACES AND PROJECTIVE SPACES}

\subsection{DEFINITION OF AFFINE SPACES}

DEFINITION 1. Given a left (resp. right) vector space T over a field K, an affine space attached to T is any homogeneous space E of the additive group T (I, §5, no. 5) such that 0 is the only operator in T leaving invariant all the elements of E (that is, T operates faithfully and transitively on E). Under these conditions, T is called the translation space of E and its elements are called the translations of E (or free vectors of E).

In what follows we shall confine our attention to the case where T is a left vector space over K. The dimension (over K) of the translation vector space T of an affine space E is called the dimension of E (over K) and is denoted by dim E or dim \(_{K}\) E. An affine space of dimension one (resp. two) is called an affine line (resp. an affine plane). The elements of an affine space are also called points.

Under the conditions of Definition 1, for \(\mathbf{t} \in \mathbf{T}\) and \(a \in \mathbf{E}\) we shall denote by \(\mathbf{t} + a\) or \(a + \mathbf{t}\) the transform of the point \(a\) under \(\mathbf{t}\) . Then the relations

\[
\mathbf {s} + (\mathbf {t} + a) = (\mathbf {s} + \mathbf {t}) + a, \quad 0 + a = a\tag{1}
\]

hold for \(\mathbf{s} \in \mathrm{T}\) , \(\mathbf{t} \in \mathrm{T}\) , \(a \in \mathrm{E}\) . The mapping \(x \mapsto x + \mathbf{t}\) is a bijection of \(\mathbf{E}\) onto itself, which we identify with \(\mathbf{t}\) . Definition 1 moreover implies that, for all \(a \in \mathrm{E}\) , the mapping \(\mathbf{t} \mapsto \mathbf{t} + a\) is a bijection of \(\mathrm{T}\) onto \(\mathbf{E}\) . In other words, given two points \(a\) , \(b\) of \(\mathbf{E}\) , there exists one and only one translation \(\mathbf{t}\) such that \(b = \mathbf{t} + a\) ; we shall denote this translation by \(b - a\) ; then the formulae

\[
a - a = 0, \quad a - b = - (b - a), \quad b = (b - a) + a\tag{2}
\]

\[
(c - b) + (b - a) = c - a
\]

hold for \(a \in \mathbf{E}\) , \(b \in \mathbf{E}\) , \(c \in \mathbf{E}\) . If four points \(a, b, a', b'\) of \(\mathbf{E}\) are such that \(b - a = b' - a'\) , the formula

\[
b ^ {\prime} = \left(b ^ {\prime} - b\right) + (b - a) + a = \left(b ^ {\prime} - a ^ {\prime}\right) + \left(a ^ {\prime} - a\right) + a
\]

and the commutativity of addition in T show that \(b' - b = a' - a\) .

Given a point \(a \in E\) , the mapping \(x \mapsto x - a\) is a bijection of \(E\) onto \(T\) ; when \(E\) is identified with \(T\) under this mapping, we say that \(E\) is considered as the vector space obtained by taking \(a\) as origin in \(E\) . Conversely, every vector space \(T\) has canonically the structure of an affine space attached to \(T\) , namely the homogeneous space structure corresponding to the subgroup \(\{0\}\) of \(T\) (I, § 5, no. 6).

Remark. The definitions of this no. and some of the results which follow extend immediately to the case where, instead of a vector space T, we consider an arbitrary commutative group with operators T.

\subsection{BARYCENTRIC CALCULUS}

PROPOSITION 1. Let \((x_{t})_{t\in I}\) be a family of points in an affine space \(\mathbf{E}\) and \((\lambda_t)_{t\in I}\) a family of elements of \(\mathbf{K}\) of finite support such that \(\sum_{i\in I}\lambda_i = 1\) (resp. \(\sum_{i\in I}\lambda_i = 0\) ). If \(a\) is any point of \(\mathbf{E}\) , the point \(x\in \mathbf{E}\) defined by

\[
x - a = \sum_ {i \in I} \lambda_ {i} (x _ {i} - a)
\]

(resp. the free vector \(\sum_{i\in I}\lambda_i(x_i - a))\) is independent of the point considered.

If \(a'\) is another point of \(E\) , then

\[
\sum_ {i} \lambda_ {i} (x _ {i} - a ^ {\prime}) = \sum_ {i} \lambda_ {i} ((x _ {i} - a) + (a - a ^ {\prime})) = \sum_ {i} \lambda_ {i} (x - a) + \left(\sum_ {i} \lambda_ {i}\right) (a - a ^ {\prime}).
\]

If \(\sum_{i}\lambda_{i} = 1\) , then \(\sum_{i}\lambda_{i}(x_{i} - a) = (x - a) + (a - a') = x - a'\) ; if \(\sum_{i}\lambda_{i} = 0\) , then \(\sum_{i}\lambda_{i}(x - a') = \sum_{i}\lambda_{i}(x - a)\) ; whence the proposition.

Under the conditions of Proposition 1, the point x defined by

\[
x - a = \sum_ {i \in I} \lambda_ {i} (x _ {i} - a)
\]

\(\left(\text{resp. the free vector} \sum_{t \in I} \lambda_t(x_t - a)\right)\) is denoted by \(\sum_{t \in I} \lambda_t x_t\) .

Thus in particular the notation \(b - a\) introduced in no. 1 is recovered. When \(\sum_{i} \lambda_{i} = 1\) , the point \(x = \sum_{i} \lambda_{i} x_{i}\) is called the barycentre of the points \(x_{i}\) given the masses \(\lambda_{i}\) .

Given \(m\) points \(a_1, \ldots, a_m\) of \(\mathbf{E}\) , whose number \(m\) is not a multiple of the characteristic of \(\mathbf{K}\) (V, § 1), the point \(g = \sum_{i=1}^{m} \frac{1}{m} a_i\) is called (by an abuse of language) the barycentre of the points \(a_i\) ( \(1 \leqslant i \leqslant m\) ) (for \(m = 2\) , we say ``midpoint'' instead of ``barycentre''); it is characterized by the relation

\[
\sum_ {i = 1} ^ {m} \left(a _ {i} - g\right) = 0.
\]

\subsection{AFFINE LINEAR VARIETIES}

DEFINITION 2. Given an affine space E, a subset V of E is called an affine linear variety (or simply a linear variety or an affine subset of E) if, for every family \((x_{t})_{t\in I}\) of points of V and every family \((\lambda_t)_{t\in I}\) of elements of K of finite support such that \(\sum_{t\in I}\lambda_t = 1\) , the barycentre \(\sum_{t\in I}\lambda_t x_t\) belongs to V.

It amounts to the same to say that the condition of Definition 2 holds for every finite family of points of V.

The empty set is a linear variety; every intersection of linear varieties is a linear variety.

Let V be a non-empty subset of E and \(a\) a point of V; the relation

\[
x - a = \sum_ {i = 1} ^ {n} \lambda_ {i} (x _ {i} - a)
\]

means that \(x\) is a barycentre \(\sum_{i=1}^{n} \lambda_i x_i + \left(1 - \sum_{i=1}^{n} \lambda_i\right)a\) of the family consisting of the \(x_i\) and \(a\) . Therefore:

PROPOSITION 2. For a non-empty subset V of an affine space E to be a linear variety, it is necessary and sufficient that V be a vector subspace for the vector space structure on E obtained by taking a point of V as origin.

In particular, the non-empty affine linear varieties of a vector space T (considered as an affine space) are just the translates of the vector subspaces of T; the vector subspaces of T are therefore the linear varieties containing 0.

Let V be a non-empty linear variety of the affine space E; the set of free vectors x - y, where x and y run through V, is a vector subspace D of the translation space T of E called the direction of V: for, if \(a \in V\) , then

\[
x - y = (x - a) - (y - a)
\]

and our assertion follows from Proposition 2. It is immediate that D operates faithfully and transitively on V, which therefore has canonically the structure of an affine space attached to D. By the dimension of the affine variety V, we mean the dimension of V with this affine space structure, that is the dimension of the vector space D. The linear varieties of dimension 0 are the points of E; those of dimension 1 (resp. 2) are called lines (resp. planes) of E.

Every vector \(\neq0\) belonging to the direction of a line is called a direction vector of this line; its components with respect to a basis of T form what is called a system of direction parameters of the line in question.

The codimension of a linear variety V in E is the codimension of its direction D in T; a linear variety of codimension 1 in E is called an (affine) hyperplane of E.

Two linear varieties with the same direction are called parallel; it amounts to the same to say that one is derived from the other by translation. If V is a linear variety in T (considered as an affine space), its direction is the linear variety parallel to V and containing 0.

PROPOSITION 3. Given a family \((a_{t})_{t\in I}\) of points of an affine space \(\mathbf{E}\) , the set \(\mathbf{V}\) of barycentres \(\sum_{i\in I}\lambda_i a_i\) (( \(\lambda_i\) ) of finite support, \(\sum_{i\in I}\lambda_i = 1\) ) is a linear variety of \(\mathbf{E}\) .

If the family \((a_{t})\) is empty, then \(\mathbf{V} = \varnothing\) because of the condition \(\sum_{i}\lambda_{i} = 1\) . It may therefore be assumed that the family \((a_{t})\) is non-empty and in this case the proposition is obvious, taking one of the \(a_{t}\) as origin in E.

The variety V is obviously the smallest linear variety containing the \(a_{t}\) ; it is said to be generated by the family \((a_{t})\) and this family is called a generating system of V.

In the notation of Proposition 3, assuming the family \((a_{\iota})\) is non-empty, for the expression for every point \(x \in V\) in the form \(x = \sum_{\iota} \lambda_{\iota} a_{\iota}\) to be unique, it is necessary and sufficient that, denoting an arbitrary index of I by \(\kappa\) , the family of vectors \(a_{\iota} - a_{\kappa}\) , where \(\iota\) runs through the set of indices \(\neq \kappa\) , be free in T. Then the family \((a_{\iota})_{\iota \in I}\) of points of E is said to be affinely free (or that its elements form an affinely free system, or are affinely independent) and that \(\lambda_{\iota}\) is the barycentric coordinate of \(x\) of index \(\iota\) with respect to the affinely free family \((a_{\iota})\) .

A family \((a_{i})_{i\in I}\) of points of E which is not affinely free is said to be affinely related.

PROPOSITION 4. For a non-empty family \((a_{t})_{t\in I}\) of points in an affine space \(E\) to be affinely related, it is necessary and sufficient that there exist a family \((\lambda_t)_{t\in I}\) of elements not all zero in \(K\) , of finite support, such that \(\sum_{i\in I}\lambda_i = 0\) and \(\sum_{i\in I}\lambda_i a_i = 0\) .

Given an index \(\kappa \in I\) , to say that the family of vectors \((a_{t} - a_{\kappa})\) , where \(t\) runs through the set of indices \(\neq \kappa\) , is related in \(T\) , means that there exists a family of scalars \((\lambda_{t})_{t \neq \kappa}\) not all zero such that \(\sum_{t \neq \kappa} \lambda_{t}(a_{t} - a_{\kappa}) = 0\) , which may also be written \(\sum_{t \in I} \lambda_{t} a_{t} = 0\) , with \(\lambda_{\kappa} = -\sum_{t \neq \kappa} \lambda_{t}\) , in other words \(\sum_{t \in I} \lambda_{t} = 0\) .

PROPOSITION 5. For a non-empty family \((a_{t})_{t\in I}\) of points of an affine space \(E\) to be affinely free, it is necessary and sufficient that, for every index \(\kappa\in I\) , \(a_{\kappa}\) do not belong to the linear variety generated by the \(a_{t}\) of index \(\neq\kappa\) .

The proposition is obvious if I has only a single element. Otherwise, taking as origin in E one of the \(a_{t}\) of index \(\neq\kappa\) , the proposition follows from §7, no.1, Remark.

\subsection{AFFINE LINEAR MAPPINGS}

DEFINITION 3. Given two affine spaces E, E' attached to two vector spaces T, T' over the same field K, a mapping u of E into E' is called an affine linear mapping (or an affine mapping) if, for every family \((x_{t})_{t\in I}\) of points of E and every family \((\lambda_{t})_{t\in I}\) such that \(\sum_{i\in I}\lambda_{i}=1\) ,

\[
u \left(\sum_ {i \in I} \lambda_ {i} x _ {i}\right) = \sum_ {i \in I} \lambda_ {i} u (x _ {i}).\tag{3}
\]

PROPOSITION 6. Let \(u\) be an affine mapping of \(E\) into \(E'\) . There exists one and only one linear mapping \(v\) of \(T\) into \(T'\) such that

\[
u (x + \mathbf {t}) = u (x) + v (\mathbf {t})
\]

for all \(x \in \mathbf{E}\) , \(\mathbf{t} \in \mathbf{T}\) .

Let \(a\) be any point of E. The mapping

\[
\mathbf {t} \mapsto u (a + \mathbf {t}) - u (a)
\]

is a linear mapping \(v_{a}\) of T into T', for we may write

\[
\begin{array}{c} a + \lambda \mathbf {t} = \lambda (a + \mathbf {t}) + (1 - \lambda) a \\ a + \mathbf {s} + \mathbf {t} = (a + \mathbf {s}) + (a + \mathbf {t}) - a \end{array}
\]

and it follows from (3) that \(v_{a}(\lambda \mathbf{t}) = \lambda v_{a}(\mathbf{t})\) and \(v_{a}(\mathbf{s} + \mathbf{t}) = v_{a}(\mathbf{s}) + v_{a}(\mathbf{t})\) . Moreover, if \(b\) is another point of \(\mathbf{E}\) , then \(v_{a} = v_{b}\) ; for the relation

\[
(a + \mathbf {t}) - a + b = b + \mathbf {t}
\]

implies

\[
u (a + \mathbf {t}) - u (a) + u (b) = u (b + \mathbf {t})
\]

that is \(u(a + \mathbf{t}) - u(a) = u(b + \mathbf{t}) - u(b)\) . Whence the existence of \(v\) ; the uniqueness is immediate.

v is called the linear mapping of T into T' associated with u. Conversely, for every linear mapping v of T into T' and every ordered pair of points \(a \in E\) , \(a' \in E'\) , it is immediately verified that

\[
x \mapsto a ^ {\prime} + v (x - a)
\]

is an affine mapping of E into \(E'\) whose associated linear mapping is v. To say that u is an affine mapping of E into \(E'\) therefore also means that, if an arbitrary point a in E and the point \(u(a)\) in \(E'\) are taken as origins, u is a linear mapping for the two vector spaces thus obtained.

Let \(\mathbf{E}''\) be a third affine space, \(\mathbf{T}''\) its translation space, \(u'\) an affine mapping of \(\mathbf{E}'\) into \(\mathbf{E}''\) and \(v'\) the linear mapping of \(\mathbf{T}'\) into \(\mathbf{T}''\) associated with \(u'\) . Clearly \(u' \circ u\) is an affine mapping of \(\mathbf{E}\) into \(\mathbf{E}''\) ; moreover, for \(a \in \mathbf{E}\) and \(\mathbf{t} \in \mathbf{T}\) ,

\[
u ^ {\prime} (u (a + \mathbf {t})) = u ^ {\prime} (u (a) + v (\mathbf {t})) = u ^ {\prime} (u (a)) + v ^ {\prime} (v (\mathbf {t}))
\]

and hence \(v'\circ v\) is the linear mapping of T into \(T''\) associated with \(u'\circ u\) . For an affine mapping u to be bijective, it is necessary and sufficient that the associated linear mapping v be so, and \(u^{-1}\) is then an affine mapping whose associated linear mapping is \(v^{-1}\) .

In particular, the affine bijections of E onto itself form a group G, called the affine group of E. The mapping which associates with \(u \in G\) the linear mapping v associated with u is, by the above, a homomorphism of G onto the linear group \(\mathbf{GL}(\mathrm{T})\) . If u is a translation, v is the identity and conversely. Hence, the kernel of the above homomorphism is the translation group T of E which is therefore a normal subgroup of G.

If \(u \in G\) , the automorphism \(t \mapsto utu^{-1}\) of T (where t is identified with the translation \(x \mapsto x + t\) ) is the linear mapping v associated with u. For \(x \in E\) and \(t \in T\) , by definition

\[
x + u \mathbf {t} u ^ {- 1} = u (u ^ {- 1} (x) + \mathbf {t}) = u (u ^ {- 1} (x)) + v (\mathbf {t}) = x + v (\mathbf {t})
\]

and hence \(utu^{-1} = v(t)\) .

Let \(a \in E\) and \(G_a\) be the subgroup of \(G\) consisting of the \(u \in G\) such that \(u(a) = a\) . If \(E\) is identified with \(T\) by taking \(a\) as origin, \(G_a\) is identified with \(\mathbf{GL}(T)\) . Every \(u \in G\) can be expressed uniquely in the form \(u = t_1 u_1\) (resp. in the form \(u = u_2 t_2\) ), where \(u_1, u_2\) are in \(G_a\) and \(t_1, t_2\) in \(T\) : for, writing \(t_1 = u(a) - a\) , \(u^{-1} t_1 \in G_a\) , whence the existence of \(u_1\) and \(t_1\) ; the existence of \(u_2\) and \(t_2\) is obtained analogously. The uniqueness follows from the fact that \(G_a \cap T\) reduces to the identity element of \(G\) . Moreover

\[
\mathbf {t} _ {1} u _ {1} = u _ {1} (u _ {1} ^ {- 1} \mathbf {t} _ {1} u _ {1})
\]

whence \(u_{2}=u_{1}\) , \(t_{2}=u_{1}^{-1}t_{1}u_{1}\) . Finally, the linear mappings associated with u and \(u_{1}\) are the same and hence, if as above \(G_{a}\) is identified with \(\mathbf{GL}(T)\) , \(u_{1}\) is the linear mapping from T to itself associated with u. It is thus seen that G is the semi-direct product of \(G_{a}\) by T (I, §6, no. 1).

Let E, E' be two affine spaces over K. The direct (resp. inverse) image of a linear variety of E (resp. E') under an affine mapping u of E into E' is a linear variety of E' (resp. E); the rank of u is by definition the dimension of \(u(\mathrm{E})\) ; it is equal to the rank of the linear mapping associated with u. If V, V' are linear varieties of the same finite dimension m in E, E' respectively, there exists an affine mapping \(u\) of \(\mathbf{E}\) into \(\mathbf{E}'\) such that \(u(\mathbf{V}) = \mathbf{V}'\) : taking as origins in \(\mathbf{E}\) and \(\mathbf{E}'\) points of \(\mathbf{V}\) and \(\mathbf{V}'\) respectively, then taking in \(\mathbf{E}\) (resp. \(\mathbf{E}'\) ) a basis whose first \(m\) vectors form a basis of \(\mathbf{V}\) (resp. \(\mathbf{V}'\) ), the proposition follows immediately from § 1, no. 11, Corollary 3 to Proposition 17.

As the field K has canonically a left vector space structure (of dimension 1) over K, it can be considered as an affine space of dimension 1. An affine mapping of an affine space D (over K) into the affine space K is also called an affine linear function (or an affine function). If a point \(a\) is taken as origin in E, every affine function on E can then be written uniquely as \(x \mapsto \alpha + v(x)\) , where \(\alpha \in \mathbf{K}\) and \(v\) is a linear form on the vector space E thus obtained; the affine functions on E therefore form a right vector space over K of dimension \(1 + \dim \mathbf{E}\) . If \(u\) is a non-constant affine function on E and \(\lambda \in \mathbf{K}\) , the set of \(x \in \mathbf{E}\) satisfying the equation \(u(x) = \lambda\) is a hyperplane; conversely, for every

hyperplane H in E, there exists an affine function \(u_{0}\) on E such that \(\mathbf{H} = \bar{u}_{0}^{-1}(0)\) and every affine function u such that \(\mathbf{H} = \bar{u}^{-1}(0)\) is of the form \(u_{0}\mu\) , where \(\mu \in K\) (§7, no. 5, Proposition 11). If u is an affine function on E, the hyperplanes with equations \(u(x) = \alpha\) and \(u(x) = \beta\) are parallel.

\subsection{DEFINITION OF PROJECTIVE SPACES}

DEFINITION 4. Given a left (resp. right) vector space V over a field K, the left (resp. right) projective space derived from V, denoted by P(V), is the quotient of the complement V - \{0\} of \{0\} in V by the equivalence relation Δ(V) ``there exists λ ≠ 0 in K such that y = λx (resp. y = xλ) between x and y in V - \{0\}.

When \(\mathbf{V} = \mathrm{K}_s^{n + 1}\) , we also write \(\mathbf{P}_n(\mathbf{K})\) instead of \(\mathbf{P}(\mathrm{K}_s^{n + 1})\) and \(\Delta_n(\mathbf{K})\) instead of \(\Delta (\mathbf{V})\) .

Definition 4 can also be expressed by saying that \(\mathbf{P}(\mathbf{V})\) is the set of lines (passing through 0) in V with the origin removed; \(\mathbf{P}(\mathbf{V})\) is therefore canonically identified with the set of lines (passing through 0) in V. The elements of a projective space are called the points of that space.

When V is of dimension n, the integer n - 1 is called the dimension of the projective space \(\mathbf{P}(\mathbf{V})\) if n is finite, and the cardinal n otherwise; this cardinal is denoted by \(\dim_{\mathbb{K}} \mathbf{P}(\mathbf{V})\) or \(\dim \mathbf{P}(\mathbf{V})\) . Thus a projective space of dimension -1 is empty and a projective space of dimension 0 is a single point. A projective space of dimension 1 (resp. 2) is called a projective line (resp. projective plane).

Henceforth we shall only consider left projective spaces.

\subsection{HOMOGENEOUS COORDINATES}

Let V be a vector space of finite dimension \(n + 1\) over K, \(\mathbf{P}(\mathbf{V})\) the projective space of dimension n derived from V and \((e_{i})_{0 \leqslant i \leqslant n}\) a basis of V. Let \(\pi\) denote the canonical mapping of V - \{0\} onto the quotient set \(\mathbf{P}(\mathbf{V})\) . For every point \(x = \sum_{i=0}^{n} \xi_i e_i\) of \(V - \{0\}\) , \((\xi_0, \xi_1, \ldots, \xi_n)\) is called a system of homogeneous coordinates of the point \(\pi(x)\) with respect to the basis \((e_i)\) of \(V\) . Every system \((\xi_i)\) of \(n + 1\) elements not all zero of \(K\) is therefore a system of homogeneous coordinates of a point of \(P(V)\) with respect to \((e_i)\) ; for two such systems \((\xi_i), (\xi_i')\) to be systems of homogeneous coordinates of the same point of \(P(V)\) with respect to the same basis \((e_i)\) , it is necessary and sufficient that there exists an element \(\lambda \neq 0\) of \(K\) such that \(\xi_i' = \lambda \xi_i\) for \(0 \leqslant i \leqslant n\) .

This definition is immediately generalized to the case where V is infinite-dimensional.

Given another basis \((\bar{e}_i)\) of V such that \(e_i = \sum_{j=0}^{n} \alpha_{ij} \bar{e}_j\) ( \(0 \leqslant i \leqslant n\) ) and a system \((\xi_i)\) of homogeneous coordinates of \(\pi(x)\) with respect to the basis \((e_i)\) , for a system \((\bar{\xi}_i)\) of \(n+1\) elements of K to be a system of homogeneous coordinates of \(\pi(x)\) with respect to the basis \((\bar{e}_i)\) , it is necessary and sufficient that there exist \(\lambda \neq 0\) in K such that

\[
\lambda \bar {\xi} _ {i} = \sum_ {j = 0} ^ {n} \xi_ {j} \alpha_ {j i} \quad \text { for } 0 \leqslant i \leqslant n.
\]

In particular, if \(e_i = \gamma_i \bar{e}_i\) with \(\gamma_i \neq 0\) ( \(0 \leqslant i \leqslant n\) ), then \(\bar{\xi}_i = \mu \xi_i \gamma_i\) where \(\mu \neq 0\) .

\subsection{PROJECTIVE LINEAR VARIETIES}

Let W be a vector subspace of a vector space V; the canonical image of \(W - \{0\}\) in the projective space \(\mathbf{P}(V)\) derived from V is called a projective linear variety (or simply a linear variety when no confusion is to be feared); as the equivalence relation \(\Delta(W)\) on \(W - \{0\}\) is induced by the relation \(\Delta(V)\) , the projective linear variety the image of \(W - \{0\}\) in \(\mathbf{P}(V)\) can be identified with the projective space \(\mathbf{P}(W)\) derived from W and hence we may speak of the dimension of such a variety. In a projective space \(\mathbf{P}(V)\) , the canonical image of a hyperplane (with the origin removed) of V is a linear variety called a projective hyperplane (or simply a hyperplane); if \(\mathbf{P}(V)\) is of finite dimension n the hyperplanes in \(\mathbf{P}(V)\) are the linear varieties of dimension n - 1.

Every proposition concerning vector subspaces of a vector space goes over to a proposition concerning projective linear varieties. For example, if a projective space \(\mathbf{P}(\mathbf{V})\) is of finite dimension n and \((e_{i})_{0\leqslant i\leqslant n}\) is a basis of V, every linear variety \(\mathbf{L}\subset\mathbf{P}(\mathbf{V})\) of dimension r can be defined by a system of n-r homogeneous linear equations

\[
\sum_ {i = 0} ^ {n} \xi_ {i} \alpha_ {i j} = 0 \quad (1 \leqslant j \leqslant n - r)\tag{4}
\]

between the homogeneous coordinates \(\xi_{i}\) ( \(0 \leqslant i \leqslant n\) ) of a point of \(\mathbf{P}(\mathbf{V})\) with respect to the basis \((e_{i})\) , the left hand sides of (4) being linearly independent forms on V. In particular, a projective hyperplane is defined by a single homogeneous linear equation with coefficients not all zero. Conversely, the points of \(\mathbf{P}(\mathbf{V})\) satisfying an arbitrary system of homogeneous linear equations with respect to the \(\xi_{i}\) form a linear variety L; if the system considered consists of \(k \leqslant n + 1\) equations, L is of dimension \(\geqslant n - k\) .

Every intersection of linear varieties of \(\mathbf{P}(\mathbf{V})\) is a linear variety; for every subset A of \(\mathbf{P}(\mathbf{V})\) , there exists a smallest linear variety L containing A; it is called the linear variety generated by A and A is called a generating system of L.

If W is the vector subspace of V generated by \(\pi^{1}(A)\) , then L = P(W).

If \(\mathbf{L}\) and \(\mathbf{M}\) are any two linear varieties in \(\mathbf{P}(\mathbf{V})\) and \(\mathbf{N}\) the variety generated by \(\mathbf{L} \cup \mathbf{M}\) , then (§ 7, no. 3, Corollary 3 to Proposition 4)

\[
\dim \mathbf {L} + \dim \mathbf {M} = \dim (\mathbf {L} \cap \mathbf {M}) + \dim \mathbf {N}.\tag{5}
\]

In particular, if \(\mathbf{P}(\mathbf{V})\) is finite-dimensional and \(\dim \mathbf{L} + \dim \mathbf{M} \geqslant \dim \mathbf{P}(\mathbf{V})\) , it follows from (5) that \(\mathbf{L} \cap \mathbf{M}\) is non-empty.

Let \((x_{i}), (y_{i})\) be two families of points in the vector space V with the same indexing set, such that \(y_{i} = \lambda_{i}x_{i}\) , where \(\lambda_{i}^{*} \neq 0\) for all i. If the family \((x_{i})\) is free, so is \((y_{i})\) and conversely; then it is said that the family of points \(\pi(x_{i})\) of P(V) is projectively free (or simply free). It amounts to the same to say that for every index K, the point \(\pi(x_{K})\) does not belong to the linear variety generated by the \(\pi(x_{i})\) for \(i \neq K\) . A family of points of P(V) which is not projectively free is called projectively related (or simply related).

For a family \((x_{i})\) of points of \(\mathbf{V} - \{0\}\) to be such that the family \((\pi (x_{i}))\) is projectively free and generates \(\mathbf{P}(\mathbf{V})\) , it is necessary and sufficient that \((x_{i})\) be a basis of \(\mathbf{V}\) . If \(\mathbf{P}(\mathbf{V})\) is of dimension \(n\) the number of elements in such a family is therefore \(n + 1\) . Note that giving such a family \((\pi (x_{i}))\) in \(\mathbf{P}(\mathbf{V})\) does not determine (even to within a left factor) the homogeneous coordinates of a given point of \(\mathbf{P}(\mathbf{V})\) with respect to a basis \((y_{i})\) of \(\mathbf{V}\) such that \(\pi (y_{i}) = \pi (x_{i})\) for all \(\iota\) (cf.~no. 6).

\subsection{PROJECTIVE COMPLETION OF AN AFFINE SPACE}

Let V be a (left) vector space over a field K and consider the vector space \(K_{s} \times V\) over K; the projective space \(\mathbf{P}(\mathbf{K}_{s} \times \mathbf{V})\) is called the projective space canonically associated with the vector space V. If V is of dimension n, \(\mathbf{P}(\mathbf{K}_{s} \times \mathbf{V})\) is of the same dimension n.~Consider in \(K_{s} \times V\) the affine hyperplane \(V_{1} = \{1\} \times V\) , whose direction (no. 3) is the subspace \(V_{0} = \{0\} \times V\) ; if a line (passing through 0) of \(K_{s} \times V\) is not contained in \(V_{0}\) , it contains a point \((\alpha, x)\) with \(\alpha \neq 0\) and \(x \in V\) , hence it contains also the point \(\alpha^{-1}(\alpha, x) = (1, \alpha^{-1}x)\) of \(V_{1}\) ; the converse is immediate and it is seen that there is a one-to-one correspondence between the points of \(V_{1}\) and the lines (passing through 0) of \(\mathbf{K}_s \times \mathbf{V}\) not contained in \(\mathbf{V}_0\) , each of the latter meeting \(\mathbf{V}_1\) in one and only one point. It follows that the mapping \(x \mapsto \phi(x) = \pi(1, x)\) is an injection (called canonical) of \(\mathbf{V}\) into the projective space \(\mathbf{P}(\mathbf{K}_s \times \mathbf{V})\) ; \(\mathbf{V}\) is often identified with its image under this injection. The complement of \(\phi(\mathbf{V})\) in \(\mathbf{P}(\mathbf{K}_s \times \mathbf{V})\) is the projective hyperplane \(\mathbf{P}(\mathbf{V}_0)\) called the hyperplane at infinity of \(\mathbf{P}(\mathbf{K}_s \times \mathbf{V})\) (or of \(\mathbf{V}\) , by an abuse of language); its points are also called the ``points at infinity'' of \(\mathbf{P}(\mathbf{K}_s \times \mathbf{V})\) (or of \(\mathbf{V}\) ). If \((a_t)\) is a basis of \(\mathbf{V}\) and in \(\mathbf{K}_s \times \mathbf{V}\) the basis is taken consisting of the elements \(e_t = (0, a_t)\) and the element \(e_\omega = (1, 0)\) , the points at infinity in \(\mathbf{P}(\mathbf{K}_s \times \mathbf{V})\) are those whose homogeneous coordinate of index \(\omega\) is 0.

Let M be an affine linear variety in V (no. 3) and D its direction; the canonical image \(\phi(M)\) of M in \(\mathbf{P}(\mathbf{K}_s \times \mathbf{V})\) is contained in the canonical image \(\overline{\mathbf{M}} = \pi(\mathbf{M}_2)\) of the vector subspace \(\mathbf{M}_2\) of \(\mathbf{K}_s \times \mathbf{V}\) generated by the affine variety \(\mathbf{M}_1 = \{1\} \times \mathbf{M}\) of \(\mathbf{K}_s \times \mathbf{V}\) . More precisely, if \((a_t)\) is an affinely free system of M generating M, the elements \((1, a_t)\) form a basis of \(\mathbf{M}_2\) and therefore \(\overline{\mathbf{M}}\) is just the projective linear variety generated by \(\phi(M)\) ; if M is finite-dimensional, \(\overline{\mathbf{M}}\) has the same dimension as M. The complement of \(\phi(M)\) in \(\overline{\mathbf{M}}\) is the intersection of \(\overline{\mathbf{M}}\) and the hyperplane at infinity and is equal to the canonical image \(\pi(\mathbf{M}_0)\) , where \(\mathbf{M}_0 = \{0\} \times \mathbf{D}\) .

Conversely, let N be a projective linear variety not contained in the hyperplane at infinity and let \(\mathbf{R} = \overline{\pi}^{-1}(\mathbf{N})\) ; \(R \cap V_{1}\) is an affine linear variety of \(K \times V\) of the form \(\{1\} \times M\) , where M is an affine linear variety of V and it is immediately seen that N is the affine linear variety \(\overline{M}\) generated by \(\phi(M)\) .

There is therefore a one-to-one correspondence between the affine linear varieties of V and the projective linear varieties of \(\mathbf{P}(\mathbf{K}_{s} \times \mathbf{V})\) not contained in the hyperplane at infinity; for two affine linear varieties of V to be parallel, it is necessary and sufficient that the projective linear varieties which they generate have the same intersection with the hyperplane at infinity (which is sometimes expressed by saying that the two affine linear varieties in question have the same points at infinity).

\subsection{EXTENSION OF RATIONAL FUNCTIONS}

If the results of no. 8 are applied to the vector space \(\mathbf{V} = \mathbf{K}_s\) of dimension 1, it is seen that there exists a canonical injection \(\phi\) of \(\mathbf{K}_s\) into the projective line \(\mathbf{P}_1(\mathbf{K}) = \mathbf{P}(\mathbf{K}_s \times \mathbf{K}_s)\) ; for all \(\xi \in \mathbf{K}\) , \(\phi(\xi)\) is the point with homogeneous coordinates \((1, \xi)\) with respect to the canonical basis ( \(\S 1\) , no. 11) of \(\mathbf{K}_s \times \mathbf{K}_s\) . The complement of \(\phi(\mathbf{K})\) in \(\mathbf{P}_1(\mathbf{K})\) consists of the single point with homogeneous coordinates \((0, 1)\) with respect to the above basis; it is called the ``point at infinity''. \(\mathbf{P}_1(\mathbf{K})\) is also called the projective field associated with \(\mathbf{K}\) and denoted by \(\tilde{\mathbf{K}}\) , the point at infinity in \(\tilde{\mathbf{K}}\) being denoted by \(\infty\) .

Consider in particular the case where K is a commutative field and let \(f \in \mathbf{K}(\mathbf{X})\) be a rational function in one indeterminate over \(\mathbf{K}\) (IV, §4); if \(f \neq 0\) , there is a unique expression \(f = \alpha p/q\) , where \(\alpha \in \mathbf{K}^*\) and \(p\) and \(q\) are two relatively prime monic polynomials (VII, §1); let \(m\) and \(n\) be their respective degrees and let \(r = \sup(m, n)\) . We write

\[
p _ {1} (\mathrm{T}, \mathrm{X}) = \mathrm{T} ^ {r} p (\mathrm{X} / \mathrm{T}), \quad q _ {1} (\mathrm{T}, \mathrm{X}) = \mathrm{T} ^ {r} q (\mathrm{X} / \mathrm{T});
\]

\(p_{1}\) and \(q_{1}\) are two homogeneous polynomials of degree r over K such that \(p(\mathbf{X}) = p_{1}(1, \mathbf{X})\) , \(q(\mathbf{X}) = q_{1}(1, \mathbf{X})\) . Hence, for every element \(\xi \in K\) which is not a zero of \(q(\mathbf{X})\) , \(f(\xi) = \alpha p(\xi)/q(\xi)\) is defined and we may write

\[
f (\xi) = \alpha p _ {1} (1, \xi) / q _ {1} (1, \xi) = \alpha p _ {1} (\lambda , \lambda \xi) / q _ {1} (\lambda , \lambda \xi)
\]

for all \(\lambda \neq 0\) in K. Consider then the mapping

\[
(\eta , \xi) \mapsto (q _ {1} (\eta , \xi), p _ {1} (\eta , \xi))
\]

of \(K^{2}\) into itself; it is compatible with the equivalence relation \(\Delta(\mathbf{K}^{2})\) and therefore defines, when passing to the quotients, a mapping \(\tilde{f}\) of \(\tilde{K}\) into itself which coincides with \(\xi\mapsto f(\xi)\) at the points where this rational function is defined; it is said, by an abuse of language, that \(\tilde{f}\) is the canonical extension of f to \(\tilde{K}\) .

For example, if \(f = 1 / \mathbf{X}\) , then \(\tilde{f}(0) = \infty\) and \(\tilde{f}(\infty) = 0\) ; if

\[
f = (a \mathbf {X} + b) / (c \mathbf {X} + d)
\]

with \(ad - bc \neq 0\) , then \(\tilde{f}(-d/c) = \infty\) , \(\tilde{f}(\infty) = a/c\) if \(c \neq 0\) , \(\tilde{f}(\infty) = \infty\) if \(c = 0\) . If \(f = a_0\mathbf{X}^n + \cdots + a_n\) is a polynomial of degree \(n > 0\) , then \(\tilde{f}(\infty) = \infty\) .

\subsection{PROJECTIVE LINEAR MAPPINGS}

Let V, \(V'\) be two left vector spaces over a field K, f a linear mapping of V into \(V'\) and \(N = \bar{f}^{-1}(0)\) its kernel. It is immediate that the image under f of a line (passing through 0) in V not contained in N is a line (passing through 0) in \(V'\) ; hence, on passing to the quotients, f defines a mapping g of \(\mathbf{P}(\mathbf{V}) - \mathbf{P}(\mathbf{N})\) into \(\mathbf{P}(\mathbf{V}')\) . Such a mapping is called a projective linear mapping (or, simply, a projective mapping); although it is defined on \(\mathbf{P}(\mathbf{V}) - \mathbf{P}(\mathbf{N})\) and not on \(\mathbf{P}(\mathbf{V})\) (when \(N \neq \{0\}\) ), we shall say by an abuse of language that g is a projective mapping of \(\mathbf{P}(\mathbf{V})\) into \(\mathbf{P}(\mathbf{V}')\) . The projective linear variety \(\mathbf{P}(\mathbf{N})\) , where g is not defined, is called the centre of g.

Note that, when \(g\) is defined on the whole of \(\mathbf{P}(\mathbf{V})\) (that is when \(\mathbf{N} = \{0\}\) ), \(g\) is an injection of \(\mathbf{P}(\mathbf{V})\) into \(\mathbf{P}(\mathbf{V}')\) .

When bases \((a_{\lambda})_{\lambda\in\mathbf{L}}, (b_{\mu})_{\mu\in\mathbf{M}}\) are given in V and \(V'\) respectively, a projective mapping of \(\mathbf{P}(\mathbf{V})\) into \(\mathbf{P}(\mathbf{V}')\) maps a point of \(\mathbf{P}(\mathbf{V})\) with homogeneous coordinates \(\xi_{\lambda}\) ( \(\lambda \in \mathbf{L}\) ) to a point of \(\mathbf{P}(\mathbf{V}')\) with a system of homogeneous coordinates \(\eta_{\mu}\) ( \(\mu \in \mathbf{M}\) ) of the form

\[
\eta_ {\mu} = \sum_ {\lambda \in L} \xi_ {\lambda} \alpha_ {\lambda \mu} \quad (\alpha_ {\lambda \mu} \in K).\tag{6}
\]

The centre of \(g\) is the linear variety defined by the equations

\[
\sum_ {\lambda \in L} \xi_ {\lambda} \alpha_ {\lambda \mu} = 0 \quad (\mu \in M).
\]

If C is the centre of g and M is a linear variety of \(\mathbf{P}(\mathbf{V})\) , the image under g of \(\mathbf{M} - (\mathbf{M} \cap \mathbf{C})\) is a linear variety of \(\mathbf{P}(\mathbf{V}')\) denoted (by an abuse of language) by \(g(\mathbf{M})\) . Then

\[
\dim g (\mathbf {M}) + \dim (\mathbf {M} \cap \mathbf {C}) + 1 = \dim \mathbf {M}\tag{7}
\]

(§ 7, no. 4, formula (12)). If \(\mathbf{M}'\) is a linear variety of \(\mathbf{P}(\mathbf{V}')\) , \(\bar{g}^{-1}(\mathbf{M}') \cup \mathbf{C}\) is a linear variety of \(\mathbf{P}(\mathbf{V})\) and

\[
\dim (\bar {g} ^ {- 1} (\mathbf {M} ^ {\prime}) \cup \mathbf {C}) = \dim \mathbf {C} + \dim (\mathbf {M} ^ {\prime} \cap g (\mathbf {P} (\mathbf {V}))) + 1.\tag{8}
\]

It is said, by an abuse of language, that \(\bar{g}^{-1}(\mathbf{M}') \cup \mathbf{C}\) is the inverse image of \(\mathbf{M}'\) under \(g\) .

As the values taken by a linear mapping on a basis \((e_{t})\) of V can be chosen arbitrarily in \(V'\) , it is seen that there exists a projective linear mapping of \(\mathbf{P}(\mathbf{V})\) into \(\mathbf{P}(\mathbf{V}')\) taking arbitrary values at the points \(\pi(e_{t})\) . But (even when g is everywhere defined) giving \(g(\pi(e_{t}))\) does not determine g uniquely (Exercise 10).

The composition of two projective mappings which are bijections is a projective mapping; so is the inverse mapping of such a bijection. The bijective projective mappings of a projective space \(\mathbf{P}(\mathbf{V})\) onto itself thus form a group, called the projective group of \(\mathbf{P}(\mathbf{V})\) and denoted by \(\mathbf{PGL}(\mathbf{V})\) ; we write \(\mathbf{PGL}_{n}(\mathbf{K})\) or \(\mathbf{PGL}(n,\mathbf{K})\) instead of \(\mathbf{PGL}(\mathbf{K}_{s}^{n})\) .

Remark. In a projective space \(\mathbf{P}(\mathbf{V})\) over a field \(\mathbf{K}\) , let \(\mathbf{H} = \mathbf{P}(\mathbf{W})\) be a hyperplane. There exists a bijective linear mapping \(f\) of \(\mathbf{V}\) onto \(\mathbf{K}_s \times \mathbf{W}\) such that \(f(\mathbf{W}) = \mathbf{W}\) ; let \(g\) be the projective mapping obtained from \(f\) by passing to the quotients. It has been seen (no. 8) that the complement of \(\mathbf{P}(\mathbf{W})\) in \(\mathbf{P}(\mathbf{K}_s \times \mathbf{W})\) can be identified with an affine space whose translation space is \(\mathbf{W}\) . When \(\mathbf{P}(\mathbf{V})\) is identified with \(\mathbf{P}(\mathbf{K}_s \times \mathbf{W})\) by means of \(g\) , it is said that \(\mathbf{H}\) has been taken as hyperplane at infinity in \(\mathbf{P}(\mathbf{V})\) ; the complement of \(\mathbf{H}\) in \(\mathbf{P}(\mathbf{V})\) is then identified with an affine space whose translation space is \(\mathbf{W}\) .

\subsection{PROJECTIVE SPACE STRUCTURE}

Given a set E and a field K, a (left) projective space structure on E with respect to the field K is defined by giving a non-empty set \(\Phi\) of bijections of subsets of the projective space \(\mathbf{P}(\mathbf{K}_{s}^{\mathrm{(E)}})\) onto E satisfying the following axioms:

\((\mathbf{EP}_{\mathrm{I}})\) The set of definition of every mapping \(f \in \Phi\) is a linear variety of \(\mathbf{P}(\mathbf{K}_s^{(\mathbf{E})})\) .

\((\mathrm{EP}_{\Pi})\) For every ordered pair of elements \(f, g\) of \(\Phi\) defined respectively on the linear varieties \(\mathbf{P}(\mathbf{V})\) and \(\mathbf{P}(\mathbf{W})\) , the bijection \(h = \overline{g}^{-1} \circ f\) of \(\mathbf{P}(\mathbf{V})\) onto \(\mathbf{P}(\mathbf{W})\) is a projective mapping.

\((\mathbf{EP}_{\mathrm{III}})\) Conversely, if \(f \in \Phi\) is defined on the linear variety \(\mathbf{P}(\mathbf{V})\) and \(h\) is a bijective projective mapping of \(\mathbf{P}(\mathbf{V})\) onto a linear variety \(\mathbf{P}(\mathbf{W}) \subset \mathbf{P}(\mathbf{K}_s^{\mathbf{(E)}})\) , then \(f \circ h^{-1} \in \Phi\) .

Let E be a set, \((\mathbf{V}_{\lambda})_{\lambda \in L}\) a family of vector spaces over K and suppose given for each \(\lambda \in L\) a bijection \(f_{\lambda}\) of \(\mathbf{P}(\mathbf{V}_{\lambda})\) onto E such that, for every ordered pair of indices \(\lambda, \mu, f_{\lambda} \circ f_{\mu}\) is a projective mapping of \(\mathbf{P}(\mathbf{V}_{\mu})\) onto \(\mathbf{P}(\mathbf{V}_{\lambda})\) . Then we can define on E a projective space structure with respect to K as follows: let \((e_{t})_{t \in I}\) be a basis of a space \(V_{\lambda}\) and write \(a_{t} = f_{\lambda}(\pi(e_{t}))\) ; let \(b_{t}\) be the element of index \(a_{t}\) in the canonical basis of \(\mathbf{K}_{s}^{(\mathbf{E})}\) (§1, no. 11). The relation \(t \neq k\) implies \(b_{t} \neq b_{k}\) because of the hypothesis that \(f_{\lambda}\) is bijective; hence the \(b_{t}\) form a basis of a vector subspace \(W_{0}\) of \(K_{s}^{(\mathbf{E})}\) and there therefore exists a bijective projective mapping h of \(\mathbf{P}(\mathbf{W}_{0})\) onto \(\mathbf{P}(\mathbf{V}_{\lambda})\) such that \(h(\pi(b_{t})) = \pi(e_{t})\) for all \(t \in I\) . If \(\Phi\) is taken to be the set of all bijective projective mappings \(f_{\lambda} \circ h \circ g^{-1}\) , where g runs through the set of all bijective projective mappings \(\mathbf{P}(\mathbf{W}) \subset \mathbf{P}(\mathbf{K}_{s}^{(\mathbf{E})})\) , it is immediately verified that \(\Phi\) satisfies axioms \((\mathrm{EP}_{\mathrm{I}}), (\mathrm{EP}_{\mathrm{II}})\) and \((\mathrm{EP}_{\mathrm{III}})\) . It is moreover immediate that \(\Phi\) depends neither on the choice of index \(\lambda \in L\) , nor on the choice of basis \((e_{t})\) in \(V_{\lambda}\) , nor on the choice of h.

In particular (taking L to consist of a single element), every projective space \(\mathbf{P}(\mathbf{V})\) derived from a vector space V (no. 5, Definition 4) thus has a well determined ``projective space structure'' in the sense of the definition given in this no. Hence any set with a projective space structure can be called a projective space.

\[
f \in \Phi
\]

\[
\mathbf {P} (\mathbf {V}) \subset \mathbf {P} \left(\mathrm{K} _ {\mathrm{s}} ^ {(\mathrm{E})}\right)
\]

\(f(\mathbf{M})\) is a linear variety in \(\mathbf{P}(\mathbf{V})\) in the sense of no. 7 (this property then holds for all \(f \in \Phi\) ). It follows from the above that every linear variety in a projective space has canonically a projective space structure.

A projective space \(\mathbf{E}\) is said to be of dimension \(n\) if, for all \(f\in\Phi,f(\mathbf{E})\) is a linear variety of dimension \(n\) (it suffices that this hold for one mapping \(f\in\Phi\) ).

\section{MATRICES}

\subsection{DEFINITION OF MATRICES}

DEFINITION 1. Let I, K, H be three sets; a matrix of type (I, K) with elements in H (or a matrix of type (I, K) over H) is any family \(M = (m_{\iota\kappa})_{(\iota, \kappa) \in I \times K}\) of elements of H whose indexing set is the product \(I \times K\) . For all \(\iota \in I\) , the family \((m_{\iota\kappa})_{\kappa \in K}\) is called the row of \(M\) of index \(\iota\) ; for all \(\kappa \in K\) , the family \((m_{\iota\kappa})_{\iota \in I}\) is called the column of \(M\) of index \(\kappa\) .

If I (resp. K) is finite, M is said to be a matrix with a finite number of rows (resp. columns). The set of matrices of type (I, K) over H is identified with the product \(H^{I \times K}\) .

The names ``row'' and ``column'' arise from the fact that, in the case where I and K are intervals \([1,p]\) , \([1,q]\) of N, the elements of the matrix are envisaged as set out in a rectangular array with p rows (arranged horizontally) and q columns (arranged vertically):

\[
\left( \begin{array}{c c c c c} m _ {1 1} & m _ {1 2} & \ldots & m _ {1 q} \\ m _ {2 1} & m _ {2 2} & \ldots & m _ {2 q} \\ \cdot & \cdot & \cdot & \cdot & \cdot \\ m _ {p 1} & m _ {p 2} & \ldots & m _ {p q} \end{array} \right)
\]

When p and q are explicit integers sufficiently small for it to be practicable, it is a convention that the above array is a symbol effectively denoting the matrix in question; this notation enables us to dispense with the use of indices, it being understood that the indices of an element are determined by its place in the array; for example, when we speak of the matrix

\[
\left( \begin{array}{c c c} a & b & c \\ d & e & f \end{array} \right)
\]

we mean the matrix \((m_{ij})_{1\leqslant i\leqslant 2,1\leqslant j\leqslant 3}\) such that

\[
m _ {1 1} = a, m _ {1 2} = b, m _ {1 3} = c, m _ {2 1} = d, m _ {2 2} = e, m _ {2 3} = f.
\]

Instead of matrix of type \(([1, p], [1, q])\) , we also say matrix of type \((p, q)\) , or matrix with \(p\) rows and \(q\) columns, if no confusion arises; the set of matrices of type \((p, q)\) over H is sometimes denoted by \(\mathbf{M}_{p,q}(\mathrm{H})\) .

Every matrix over H for which one of the indexing sets I, K is empty is identical with the empty family of elements of H; it is also called the empty matrix. When \(I = \{i_{0}\}\) (resp. \(K = \{k_{0}\}\) ) is a set consisting of a single element, M is called a row matrix (resp. column matrix) and the row (resp. column) index can then be suppressed in the notation; when I and K are both sets with one element, a matrix of type (I, K) is often identified with the unique element in this matrix.

A subfamily \(M' = (m_{\iota \kappa})_{(\iota, \kappa) \in J \times L}\) of a matrix \(M = (m_{\iota \kappa})_{(\iota, \kappa) \in I \times K}\) , whose indexing set is the product of a subset \(J\) of \(I\) and a subset \(L\) of \(K\) , is called a submatrix of the matrix \(M\) ; it is said to be obtained by suppressing in \(M\) the rows of index \(\iota \notin J\) and the columns of index \(\kappa \notin L\) ; conversely, \(M\) is said to be obtained by bordering \(M'\) with the rows of index \(\iota \notin J\) and the columns of index \(\kappa \notin L\) .

DEFINITION 2. The transpose of a matrix \(M = (m_{\iota \kappa})_{(\iota, \kappa) \in \mathbf{I} \times \mathbf{K}}\) , denoted by \({}^tM\) , is the matrix \((m_{\kappa \iota}')_{(\kappa, \iota) \in \mathbf{K} \times \mathbf{I}}\) over \(\mathbf{H}\) given by \(m_{\kappa \iota}' = m_{\iota \kappa}\) for all \((\iota, \kappa) \in \mathbf{I} \times \mathbf{K}\) .

It follows from this definition that the transpose of a matrix of type (I, K) is a matrix of type (K, I) and that

\[
{ } ^ { t } ( { } ^ { t } M ) = M .\tag{1}
\]

\subsection{MATRICES OVER A COMMUTATIVE GROUP}

Let G be a commutative group (written additively). The set of matrices over G, with the given indexing sets I, K, has a commutative group structure since it is the set of mappings from \(I \times K\) to G; this group is written additively, so that if \(M = (m_{\iota\kappa})\) and \(M' = (m_{\iota\kappa}')\) are two of its elements, then

\[
M + M ^ {\prime} = (m _ {\iota \kappa} + m _ {\iota \kappa} ^ {\prime});
\]

the identity element of this group is therefore the matrix all of whose elements are zero (called the zero matrix). Clearly

\[
{ } ^ { t } ( M + M ^ { \prime } ) = { } ^ { t } M + { } ^ { t } M ^ { \prime } .\tag{2}
\]

The sum of two matrices is thus only defined if the indexing sets of the rows and the columns are the same for the two matrices.

Let \(\mathbf{H}'\) , \(\mathbf{H}''\) be two sets, \(\mathbf{G}\) a commutative group (written additively) and \(f: (h', h'') \mapsto h'h''\) a mapping from \(\mathbf{H}' \times \mathbf{H}''\) to \(\mathbf{G}\) . Given two matrices

\[
M ^ {\prime} = \left(m _ {i k} ^ {\prime}\right) _ {(i, k) \in I \times K}, \quad M ^ {\prime \prime} = \left(m _ {k l} ^ {\prime \prime}\right) _ {(k, l) \in K \times L}
\]

over H' and H'' respectively such that the indexing set K of the columns of M' is finite and equal to the indexing set of the rows of M'', the product of M' and M'' via f, denoted by M'M'' or f(M', M''), is the matrix

\[
\left(\sum_ {k \in \mathbf {K}} m _ {i k} ^ {\prime} m _ {k l} ^ {\prime \prime}\right) _ {(i, l) \in \mathbf {I} \times \mathbf {L}}\tag{3}
\]

over G.

The above definition supposes that the indexing set of the columns of \(M'\) is equal to the indexing set of the rows of \(M''\) ; in particular the product \(M''M'\) has no meaning if \(I \neq L\) . In formula (3) the elements of the same

row of \(M'\) figure multiplied on the right by the elements of the same column of \(M''\) ; the multiplication is said to be made ``rows by columns''.

Let \(f^0\) be the mapping \((h'', h') \mapsto h'h''\) of \(H'' \times H'\) to G; it follows immediately from the definitions that

\[
{ } ^ { t } ( M ^ { \prime } M ^ { \prime \prime } ) = { } ^ { t } M ^ { \prime \prime } . { } ^ { t } M ^ { \prime }\tag{4}
\]

where the product on the left (resp. right) hand side is calculated via \(f\) (resp. via \(f^0\) ).

When \(\mathbf{H}'\) and \(\mathbf{H}''\) are themselves commutative groups (written additively) and \(f\) is \(\mathbf{Z}\) -bilinear ( \(\S 3\) , no. 1), the distributivity formulae

\[
\left\{ \begin{array}{l} (M ^ {\prime} + N ^ {\prime}) M ^ {\prime \prime} = M ^ {\prime} M ^ {\prime \prime} + N ^ {\prime} M ^ {\prime \prime} \\ M ^ {\prime} (M ^ {\prime \prime} + N ^ {\prime \prime}) = M ^ {\prime} M ^ {\prime \prime} + M ^ {\prime} N ^ {\prime \prime} \end{array} \right.\tag{5}
\]

are immediately verified, the indexing sets being such that the sums and products appearing are defined.

Now let \(H_{1}\) , \(H_{2}\) , \(H_{3}\) , \(H_{12}\) , \(H_{23}\) and H be commutative groups (written additively), \(f_{12}:H_{1} \times H_{2} \to H_{12}\) , \(f_{23}:H_{2} \times H_{3} \to H_{23}\) mappings and

\[
f _ {3}: \mathrm{H} _ {1 2} \times \mathrm{H} _ {3} \rightarrow \mathrm{H}, \quad f _ {1}: \mathrm{H} _ {1} \times \mathrm{H} _ {2 3} \rightarrow \mathrm{H}
\]

Z-bilinear mappings; suppose further that, for all \(x_{i} \in H_{i}\) (i = 1, 2, 3)

\[
f _ {3} (f _ {1 2} (x _ {1}, x _ {2}), x _ {3}) = f _ {1} (x _ {1}, f _ {2 3} (x _ {2}, x _ {3}))
\]

(which may also be written as above \((x_{1}x_{2})x_{3} = x_{1}(x_{2}x_{3}))\) ; then, if \(M' = (m'_{rs})\) , \(M'' = (m''_{st})\) , \(M''' = (m''_{tu})\) are matrices over \(\mathrm{H}_1\) , \(\mathrm{H}_2\) , \(\mathrm{H}_3\) respectively,

\[
(M ^ {\prime} M ^ {\prime \prime}) M ^ {\prime \prime \prime} = M ^ {\prime} (M ^ {\prime \prime} M ^ {\prime \prime \prime})\tag{6}
\]

when the products on the two sides (calculated respectively via \(f_{12}\) , \(f_{3}\) , \(f_{23}\) and \(f_{1}\) ) are defined; for

\[
\begin{array}{r l} \sum_ {t} \left(\sum_ {s} m _ {r s} ^ {\prime} m _ {s t} ^ {\prime \prime}\right) m _ {t u} ^ {\prime \prime \prime} & = \sum_ {t} \sum_ {s} (m _ {r s} ^ {\prime} m _ {s t} ^ {\prime \prime}) m _ {t u} ^ {\prime \prime \prime} = \sum_ {s} \sum_ {t} m _ {r s} ^ {\prime} (m _ {s t} ^ {\prime \prime} m _ {t u} ^ {\prime \prime \prime}) \\ & = \sum_ {s} m _ {r s} ^ {\prime} \left(\sum_ {t} m _ {s t} ^ {\prime \prime} m _ {t u} ^ {\prime \prime \prime}\right) \end{array}
\]

by virtue of the hypotheses made.

The two sides of (6) are also denoted by \(M'M''M'''\) . Analogous conventions are made for products of more than three factors.

Remark. The above formulae extend to a more general situation. To be precise:

\begin{enumerate}
\def\labelenumi{(\alph{enumi})}
\item
  Suppose \(H = \bigcup_{(\iota, \kappa) \in I \times K} G_{\iota\kappa}\) where each \(G_{\iota\kappa}\) is a commutative group written additively; then the sum \(M + M'\) may be defined when, for each ordered pair \((\iota, \kappa)\) , \(m_{\iota\kappa} \in G_{\iota\kappa}\) and \(m'_{\iota\kappa} \in G_{\iota\kappa}\) .
\item
  Let I, K, L be three sets with K assumed finite and let \(H' = \bigcup_{(i, k) \in I \times K} H'_{ik}\) , \(H'' = \bigcup_{(k, l) \in K \times L} H''_{kl}\) , \(H = \bigcup_{(i, l) \in I \times L} H_{il}\) be three sets; suppose that each \(H_{il}\) is a commutative group written additively and for each triple \((i, k, l)\) let
\end{enumerate}

\[
f _ {i k l}: \mathrm{H} _ {i k} ^ {\prime} \times \mathrm{H} _ {k l} ^ {\prime \prime} \rightarrow \mathrm{H} _ {i l}
\]

be a mapping. Then if \(M' = (m_{ik}')_{(i,k) \in I \times K}\) , \(M'' = (m_{kl}'')_{(k,l) \in K \times L}\) are matrices such that \(m_{ik}' \in H_{ik}'\) and \(m_{kl}'' \in H_{kl}''\) for all \(i, k, l\) we can define the product \(M'M''\) via the \(f_{ikl}\) . We leave to the reader the task of writing down and proving the formulae analogous to (4), (5) and (6).

\subsection{MATRICES OVER A RING}

The most important matrices in Mathematics are matrices over a ring A. The set \(A^{I \times K}\) of matrices over A corresponding to indexing sets I, K then has canonically an (A, A)-bimodule structure (§ 1, no. 14).

For every ordered pair \((i,k)\in\mathbf{I}\times\mathbf{K}\) , let \(E_{ik}\) be the matrix \((a_{jl})\) such that \(a_{ik}=1\) and \(a_{jl}=0\) for \((j,l)\neq(i,k)\) ; the \(E_{ik}\) are called the matrix units in the set of matrices \(\mathbf{A}^{\mathbf{I}\times\mathbf{K}}\) ; if I and K are finite, they form the canonical basis of this set for its left or right A-module structure (§ 1, no. 11). Clearly

\[
{ } ^ { t } E _ { i k } = E _ { k i } .
\]

Unless otherwise mentioned, the product \(M'M''\) of two matrices over A (assumed to be defined) will always be understood to be relative to the multiplication \((x,y)\mapsto xy\) in A (or, as is also said, will be ``calculated in A''). Then we have (no. 2) the associativity and distributivity formulae

\[
(X Y) Z = X (Y Z)\tag{7}
\]

\[
\left\{ \begin{array}{l} X (Y + Z) = X Y + X Z \\ (X + Y) Z = X Z + Y Z \end{array} \right.\tag{8}
\]

for three matrices X, Y, Z over A, whenever the sums and products appearing in these formulae are defined.

In particular, if \(E_{ik}\) (resp. \(E_{kl}'\) , \(E_{il}''\) ) are the matrix units in \(\mathbf{A}^{\mathbf{I} \times \mathbf{K}}\) (resp. \(\mathbf{A}^{\mathbf{K} \times \mathbf{L}}\) , \(\mathbf{A}^{\mathbf{I} \times \mathbf{L}}\) ) respectively, with \(\mathbf{I} = [1, p]\) , \(\mathbf{K} = [1, q]\) , \(\mathbf{L} = [1, r]\) , we obtain the formulae

\[
\left\{ \begin{array}{l l} E _ {i k} E _ {j l} ^ {\prime} = 0 & \text { if } k \neq j \\ E _ {i k} E _ {k l} ^ {\prime} = E _ {i l} ^ {\prime \prime}. \end{array} \right.\tag{9}
\]

Let \(\mathbf{A}^0\) be the opposite ring of \(\mathbf{A}\) and let \(a * b (= ba)\) denote the product of \(a\) and \(b\) in \(\mathbf{A}^0\) ; then, for two matrices \(X, Y\) over \(\mathbf{A}\) whose product is defined,

\[
{ } ^ { t } ( X Y ) = { } ^ { t } Y * { } ^ { t } X\tag{10}
\]

where on the right hand side \(^{t}Y\) and \(^{t}X\) are considered as matrices with elements in \(A^{0}\) ; when A is commutative, then

\[
{ } ^ { t } ( X Y ) = { } ^ { t } Y . { } ^ { t } X\tag{11}
\]

PROPOSITION 1. Let A, B be two rings and \(M = (m_{ik})_{(i,k) \in I \times K}\) and

\[
M ^ {\prime} = \left(m _ {i k} ^ {\prime}\right) _ {(i, k) \in \mathbf {I} \times \mathbf {K}}
\]

two matrices with finite indexing sets over an (A, B)-bimodule G. Suppose that for every matrix unit \(L = (a_{i})_{i \in I}\) with one row and elements in A and every matrix unit \(C = (b_{k})_{k \in K}\) with one column and elements in B, L.M.C = L.M'.C (the products being calculated via the external laws of the (A, B)-module G); then \(M = M'\) .

If \(L\) is taken to be the matrix unit \((a_s)\) with \(a_i = 1\) , \(a_s = 0\) for \(s \neq i\) , and \(C\) the matrix unit \((b_t)\) with \(b_k = 1\) , \(b_t = 0\) for \(t \neq k\) , the products \(L.M.C\) and \(L.M'.C\) are matrices with a single element respectively equal to \(m_{ik}\) and \(m'_{ik}\) .

Let A, B be two rings and \(\sigma: A \to B\) a homomorphism.

For every matrix \(M = (m_{\iota\kappa})\) over \(\mathbf{A}\) , we shall denote by \(\sigma(M)\) the matrix \((\sigma(m_{\iota\kappa}))\) over \(\mathbf{B}\) ; clearly \(\sigma(aM) = \sigma(a)\sigma(M)\) , \(\sigma(Ma) = \sigma(M)\sigma(a)\) for \(a \in \mathbf{A}\) , also \(\sigma(^t M) = ^t(\sigma(M))\) and

\[
\left\{ \begin{array}{c} \sigma (M + M ^ {\prime}) = \sigma (M) + \sigma (M ^ {\prime}) \\ \sigma (M M ^ {\prime}) = \sigma (M) \sigma (M ^ {\prime}) \end{array} \right.\tag{12}
\]

when the operations considered are defined, the products on the left and right hand sides of (12) being calculated in A and B respectively. When \(\sigma\) is denoted by \(x \mapsto x^{\sigma}\) , we write \(M^{\sigma}\) instead of \(\sigma(M)\) .

Consider in particular an anti-endomorphism \(\sigma\) of A, that is a homomorphism of A to the opposite ring \(\mathbf{A}^0\) , or a mapping of A into itself such that

\[
\sigma (a + a ^ {\prime}) = \sigma (a) + \sigma (a ^ {\prime}), \quad \sigma (a a ^ {\prime}) = \sigma (a ^ {\prime}) \sigma (a)
\]

for all \(a, a'\) in \(\mathbf{A}\) ; then, for two matrices \(M, M'\) over \(\mathbf{A}\) whose product \(MM'\) is defined,

\[
\sigma (M M ^ {\prime}) = ^ {t} (\sigma (^ {t} M ^ {\prime}). \sigma (^ {t} M))\tag{13}
\]

where the products on the two sides are calculated in A; this follows immediately from (10) and (12).

\subsection{MATRICES AND LINEAR MAPPINGS}

Let A be a ring and E a (right or left) A-module admitting a basis \((e_i)_{i\in I}\) . For every element \(x\in E\) , the matrix of \(x\) with respect to the basis \((e_i)\) , denoted by \(M(x)\) or \(\mathbf{x}\) (or sometimes simply \(x\) when no confusion can arise), is the column matrix consisting of the components \(x_i\) ( \(i\in I\) ) of \(x\) with respect to \((e_i)\) ( \(\S 1\) , no. 11); in calculations it will sometimes be convenient, in order to remember that the index \(i\) is a row index, to adjoin to it a column index taking only one value and write the matrix \(M(x)\) as \((x_{i0})\) .

We now consider two (left or right) A-modules E and F with bases \((e_{i})_{i\in I}\) and \((f_{k})_{k\in K}\) respectively; let \((f_{k}^{*})\) be the family of coordinate forms corresponding to \((f_{k})\) . For a linear mapping u of E into F, we shall define the matrix of u with respect to the bases \((e_{i})\) , \((f_{k})\) in each of the following cases:

\begin{enumerate}
\def\labelenumi{(\Alph{enumi})}
\setcounter{enumi}{3}
\item
  E and F are right A-modules, u is A-linear.
\item
  E and F are left A-modules, u is A-linear.
\end{enumerate}

In what follows, we shall attach the letter (D) (resp. (G)) to formulae applying to right (resp. left) modules.

DEFINITION 3. In each of the above two cases, the matrix of \(u\) with respect to the bases \((e_i), (f_k)\) is the matrix \(M(u) = (u_{ki})_{(k,i) \in K \times I}\) such that

\[
u _ {k i} = f _ {k} ^ {*} (u (e _ {i}))\tag{14}
\]

which is written respectively as

\[
u _ {k i} = \langle f _ {k} ^ {*}, u (e _ {i}) \rangle\tag{14 D}
\]

\[
u _ {k i} = \langle u (e _ {i}), f _ {k} ^ {*} \rangle .\tag{14 G}
\]

The column of \(M(u)\) of index i is therefore equal to \(M(u(e_{i}))\) .

Clearly if u, v are two linear mappings of E into F and \(M(u)\) , \(M(v)\) their matrices with respect to the same bases, then

\[
M (u + v) = M (u) + M (v)\tag{15}
\]

and

\[
M (\gamma u) = \gamma M (u)\tag{16}
\]

for every element \(\gamma\) of the centre \(\Gamma\) of A. In other words, once the bases \((e_i)\) , \((f_k)\) are fixed, the mapping \(u \mapsto M(u)\) is a \(\Gamma\) -module isomorphism of \(\operatorname{Hom}_{\mathbf{A}}(\mathbf{E}, \mathbf{F})\) onto a subset of the set \(\mathbf{A}^{\mathbf{K} \times \mathbf{I}}\) , equal to \(\mathbf{A}^{\mathbf{K} \times \mathbf{I}}\) if \(\mathbf{K}\) is finite.

PROPOSITION 2. Suppose I and K are finite. For every element \(x \in \mathbf{E}\) , the matrix \(M(u(x))\) with respect to the basis \((f_k)\) is given by the formula

\[
M (u (x)) = M (u). M (x)\tag{17 D}
\]

\[
{ } ^ { t } M ( u ( x ) ) = { } ^ { t } M ( x ) . { } ^ { t } M ( u ) .\tag{17 G}
\]

We verify for example (17 G). Let \(x = \sum_{i} x_{i0} e_i\) , \(u(x) = \sum_{k} y_{k0} f_k\) with \(x_{i0} \in \mathbf{A}\) , \(y_{k0} \in \mathbf{A}\) ; then \(u(x) = u\left(\sum_{i} x_{i0} e_i\right) = \sum_{i} x_{i0} u(e_i) = \sum_{i, k} x_{i0} u_{ki} f_k\) ; whence \(y_{k0} = \sum_{i} x_{i0} u_{ki}\) . In order to bring the two indices \(i\) along side one another, we consider the transpose matrices \(^t M(x) = (x'_{0i})\) , where \(x'_{0i} = x_{i0}\) and

\[
{ } ^ { t } M ( u ) = \left( u _ { i k } ^ { \prime } \right) ,
\]

where \(u_{ik}' = u_{ki}\) ; then \(y_{k0} = \sum_{i} x_{0i}' u_{ik}'\) and the right hand side is the element of index \(k\) of the matrix with one row \(^t M(x).^t M(u)\) , whence (17 G).

When A is commutative, (17 G) reduces to (17 D) by formula (4) of no. 2.

COROLLARY. Let E, F, G be three right (resp. left) modules over a ring A, \((e_{i})_{i\in I}\) , \((f_{k})_{k\in K}\) , \((g_{l})_{l\in L}\) respective finite bases of E, F, G, \(u:E\to F\) , \(v:F\to G\) two linear mappings, \(M(u)\) the matrix of u relative to the bases \((e_{i})\) , \((f_{k})\) , \(M(v)\) the matrix of v relative to the bases \((f_{k})\) , \((g_{l})\) and \(M(v\circ u)\) the matrix of \(v\circ u\) relative to the bases \((e_{i})\) , \((g_{l})\) ; then

\[
M (v \circ u) = M (v) M (u)\tag{18 D}
\]

\[
{ } ^ { t } M ( v \circ u ) = { } ^ { t } M ( u ) { } ^ { t } M ( v ) .\tag{18 G}
\]

We prove for example (18 G). For all \(x \in \mathbf{E}\) , by (17 G):

\(^{t}M(x).^{t}M(v\circ u)=^{t}M(v(u(x)))=^{t}M(u(x)).^{t}M(v)=^{t}M(x).^{t}M(u).^{t}M(v)\) by associativity; the corollary then follows from no. 3, Proposition 1 since the matrix \(^{t}M(x)\) with one row is arbitrary.

Remark (1). Formula (17 D) can be considered as a special case of (18 D). For there corresponds canonically to every \(x \in \mathbf{E}\) the linear mapping \(\theta_x: \mathbf{A}_d \to \mathbf{E}\) mapping every \(\alpha \in \mathbf{A}\) to \(x\alpha\) (§ 2, no. 1). It is immediate that the matrix \(M(\theta_x)\) with respect to the basis 1 of \(\mathbf{A}_d\) and the basis \((e_i)\) of \(\mathbf{E}\) is just the matrix \(M(x)\) ; similarly \(M(\theta_{u(x)}) = M(u(x))\) and formula (17 D) can therefore be considered as a translation of the relation

\[
\theta_ {u (x)} = u \circ \theta_ {x}.
\]

PROPOSITION 3. Let E, F be two right (resp. left) A-modules and \((e_i)_{i\in I}, (f_k)_{k\in K}\) finite bases of E and F respectively. For every linear mapping \(u\) of E into F, let \(M(u)\) be the matrix of \(u\) with respect to the bases \((e_i)\) and \((f_k)\) . Then the matrix of \(^t u: F^* \to E^*\) with respect to the dual bases \((f_k^*)\) and \((e_i^*)\) is equal to \(^t M(u)\) .

E is canonically identified with its bidual \(E^{**}\) and \((e_{i})\) with the dual basis of \((e_{i}^{*})\) ; then (supposing for example that E and F are right modules)

\[
\langle^ {t} u (f _ {k} ^ {*}), e _ {i} \rangle = \langle f _ {k} ^ {*}, u (e _ {i}) \rangle ,
\]

whence the proposition.

Remarks. (2) Let E and F be two left A-modules with bases \((e_{i})_{i\in\mathbf{I}}\) and \((f_{k})_{k\in\mathbf{K}}\) respectively. For every A-linear mapping \(u:E\to F\) , by (14 G), \(u(e_{i})=\sum_{k}u_{ki}f_{k}\) ; these relations can also be interpreted by saying that the column matrix \((u(e_{i})_{i\in\mathbf{I}})\) with elements in F is equal to the product \({}^{t}\mathbf{M}(u).(f_{k})\) , where \((f_{k})_{k\in\mathbf{K}}\) is considered as a column matrix with elements in F and the product is calculated for the mapping \(A\times F\to F\) defining the law of action on the A-module F (no. 2).

\begin{enumerate}
\def\labelenumi{(\arabic{enumi})}
\setcounter{enumi}{2}
\tightlist
\item
  Let A, B be two commutative rings and \(\sigma: A \to B\) a ring homomorphism. In the notation of Proposition 3, \((e_i \otimes 1)\) and \((f_k \otimes 1)\) are respective bases of \(E_{(B)} = E \otimes_A B\) and \(F_{(B)} = F \otimes_A B\) ( \(\S 5\) , no. 1, Proposition 4); moreover, if \((e_i^*)\) and \((f_k^*)\) are respectively the dual bases of \((e_i)\) and \((f_k)\) , then \((e_i^* \otimes 1)\) and \((f_k^* \otimes 1)\) are respectively the dual bases of \((e_i \otimes 1)\) and \((f_k \otimes 1)\) ( \(\S 5\) , no. 4). For every A-linear mapping \(u: E \to F\) , let \(M(u)\) and \(M(u \otimes 1)\) be the matrix of \(u\) with respect to \((e_i)\) and \((f_k)\) and the matrix of the B-linear mapping \(u \otimes 1\) with respect to \((e_i \otimes 1)\) and \((f_k \otimes 1)\) . It follows from \(\S 5\) , no. 4, formula (20) that
\end{enumerate}

\[
M (u \otimes 1) = \sigma (M (u)).
\]

Consider a system of a finite number of right scalar linear equations in a finite number of unknowns

\[
\sum_ {i \in I} a _ {k i} x _ {i} = b _ {k} \quad (k \in K)\tag{19}
\]

with \(a_{ki}\) , \(x_{i}\) , \(b_{k}\) in A.

Let \((e_i)_{i\in I}, (f_k)_{k\in K}\) be the canonical bases of \(\mathbf{E} = \mathbf{A}_d^{\mathrm{I}}\) and \(\mathbf{F} = \mathbf{A}_d^{\mathrm{K}}\) ; the system (19) is equivalent to the equation \(u(x) = b\) , where \(x = \sum_{i} e_i x_i\) , \(b = \sum_{k} f_k b_k\) and \(u: \mathbf{E} \to \mathbf{F}\) is the linear mapping such that the matrix \(M(u)\) with respect to the bases \((e_i)\) and \((f_k)\) are equal to \(\mathbf{A} = (a_{ki})_{(k,i)\in K\times I}\) . This matrix is called the matrix of the system of linear equations (19). Recall (§ 2, no. 8, Remarks 2 and 3), that, writing \(c_i = \sum_{k} f_k a_{ki}\) , the system (19) is equivalent to the unique vector equation

\[
\sum_ {i} c _ {i} x _ {i} = b,\tag{20}
\]

and as \(c_{i}\) is the column of index i in the matrix A, we see that to say that the system (19) admits a solution amounts to saying that the matrix \(b = (b_{k0})\) with one column is a linear combination of the columns of the matrix A.

We leave to the reader the task of formulating the analogous definitions and remarks for systems of left linear equations.

\subsection{BLOCK PRODUCTS}

The definitions of no. 4 can be generalized as follows. Let E be a (right or left) A-module, the direct sum of a family \((\mathbf{E}_i)_{i\in I}\) of submodules. For all \(x\in \mathbf{E}\) , let \(x = \sum_{i\in I}x_i\) with \(x_{i}\in \mathbf{E}_{i}\) for all \(i\in I\) ; we shall say that the column matrix \(M(x) = (x_{i})_{i\in I}\) with elements in E is the matrix of \(x\) with respect to the decomposition \((\mathbf{E}_i)_{i\in I}\) of E as a direct sum.

Let F be another A-module (E and F being both right A-modules or both left A-modules) and suppose that F is the direct sum of a family \((\mathbf{F}_{k})_{k\in\mathbb{K}}\) of submodules. For all \(u\in\mathrm{Hom}(\mathbf{E},\mathbf{F})\) and all \(x_{i}\in\mathbf{E}_{i}\) , let \(u(x_{i})=\sum_{k}u_{ki}(x_{i})\) with \(u_{ki}(x_{i})\in\mathbf{F}_{k}\) for all \(k\in K\) ; then \(u_{ki}\in\mathrm{Hom}(\mathbf{E}_{i},\mathbf{F}_{k})\) ; we shall say that the matrix \(M(u)=(u_{ki})_{(k,i)\in K\times I}\) of type (K, I) with elements in the set H the sum of the \(\mathrm{Hom}(\mathbf{E}_{i},\mathbf{F}_{k})\) is the matrix of u with respect to the decompositions \((\mathbf{E}_{i})\) and \((\mathbf{F}_{k})\) of E and F as direct sums.

With these definitions, it is obvious that if \(u, v\) are two A-linear mappings of \(E\) into \(F\) then, for matrices with respect to the same decompositions as direct sums

\[
M (u + v) = M (u) + M (v), \quad M (\gamma u) = \gamma M (u)\tag{21}
\]

for every element \(\gamma\) of the centre of A (no. 2, Remark).

Moreover, the definition of the \(u_{ki}\) shows that, if K is finite, we can write

\[
M (u (x)) = M (u). M (x)\tag{22}
\]

where \(M(u(x))\) is the matrix of \(u(x)\) with respect to the decomposition \((\mathbf{F}_k)\) , the product on the right hand side of (22) being calculated for the mappings \((t, z) \mapsto t(z)\) of \(\operatorname{Hom}(\mathbf{E}_i, \mathbf{F}_k) \times \mathbf{E}_i\) into \(\mathbf{F}_k\) (no. 2, Remark).

Let G be a third A-module, the direct sum of a family \((\mathbf{G}_{l})_{l\in L}\) of submodules, so that there corresponds to every A-linear mapping \(v:F\to G\) a matrix \(M(v)=(v_{lk})\) with respect to the decompositions \((\mathbf{F}_{k})\) and \((\mathbf{G}_{l})\) . If I, K and L are finite, then

\[
M (v \circ u) = M (v). M (u)\tag{23}
\]

where the left hand side is the matrix \((w_{li})\) of \(w = v \circ u\) with respect to the decomposition \((\mathbf{E}_i)\) and \((\mathbf{G}_l)\) and the product on the left hand side is calculated for the mappings \((t, s) \mapsto t \circ s\) of \(\operatorname{Hom}(\mathbf{F}_k, \mathbf{G}_l) \times \operatorname{Hom}(\mathbf{E}_i, \mathbf{F}_k)\) into \(\operatorname{Hom}(\mathbf{E}_i, \mathbf{G}_l)\) (no. 2, Remark). This is just formula (32) of § 1, no. 8, expressed in terms of matrices.

Finally, if I and K are assumed to be finite, \(E^{*}\) (resp. \(F^{*}\) ) is canonically identified with the direct sum of the modules \(E_{i}^{*}\) (resp. \(F_{k}^{*}\) ) ( \(\S\) 2, no. 6, Proposition 10). Then it is immediately verified that the matrix of \(^{t}u\) with respect to the decompositions \((\mathbf{F}_{k}^{*})\) and \((\mathbf{E}_{i}^{*})\) is just \((^{t}u_{kt})_{(k,t)\in\mathbb{K}\times\mathbb{I}}\) .

Suppose now that I and K are finite and further that each of the \(E_{i}\) (resp. \(F_{k}\) ) admits a finite basis. It amounts to the same to say that E (resp. F) admits a basis \((e_{r})_{r\in\mathbb{R}}\) (resp. \((f_{s})_{s\in\mathbb{S}}\) ) and that R (resp. S) admits a partition \((\mathbf{R}_{i})_{i\in\mathbb{I}}\) (resp. \((\mathbf{S}_{k})_{k\in\mathbb{K}}\) ) such that for all \(i\in I\) (resp. \(k\in K\) ), \((e_{r})_{r\in\mathbb{R}_{i}}\) is a basis of \(E_{i}\) (resp. \((f_{s})_{s\in\mathbb{S}_{k}}\) is a basis of \(F_{k}\) ). Then, if \(X=M(u)\) is the matrix of u with respect to the bases \((e_{r})_{r\in\mathbb{R}}\) and \((f_{s})_{s\in\mathbb{S}}\) , the matrix \(X_{ki}=M(u_{ki})\) with respect to the bases \((e_{r})_{r\in\mathbb{R}_{i}}\) and \((f_{s})_{s\in\mathbb{S}_{k}}\) is just the submatrix of X obtained by suppressing the rows of index \(s\notin S_{k}\) and the columns of index \(r\notin R_{i}\) . Thus we define a one-to-one correspondence

\[
X \mapsto (X _ {k i}) _ {(k, i) \in \mathbf {K} \times \mathbf {I}}\tag{24}
\]

between the set of matrices of type (S, R) with elements in A and the set of matrices of matrices \((\mathbf{X}_{ki})_{(k,i)\in\mathbf{K}\times\mathbf{I}}\) of type \(K\times I\) , where each \(X_{ki}\) is a matrix over A of type \((\mathbf{S}_{k},\mathbf{R}_{i})\) . Suppose that further G admits a finite basis \((g_{t})_{t\in\mathbf{T}}\) and that \(\mathbf{T}=(\mathbf{T}_{l})_{l\in\mathbf{L}}\) is a partition of T such that, for each \(l\in L\) , \((g_{t})_{t\in\mathbf{T}_{l}}\) is a basis of \(G_{l}\) ; let \(Y=M(v)\) be the matrix of v with respect to the bases \((f_{s})_{s\in\mathbf{S}}\) and \((g_{t})_{t\in\mathbf{T}}\) , \(Y_{lk}=M(v_{lk})\) that of \(v_{lk}\) with respect to the bases \((f_{s})_{s\in\mathbf{S}_{k}}\) and \((g_{t})_{t\in\mathbf{T}_{l}}\) , Z=M(w) the matrix of w=v∘u with respect to the bases \((e_{r})_{r\in\mathbf{R}}\) and \((g_{t})_{t\in\mathbf{T}}\) and \(Z_{li}=M(w_{li})\) that of \(w_{li}\) with respect to the bases \((e_{r})_{r\in\mathbf{R}_{i}}\) and \((g_{t})_{t\in\mathbf{T}_{l}}\) ; then it follows from (23) that the submatrices \(Z_{li}\) of Z=YX are given by

\[
Z _ {l i} = \sum_ {k} Y _ {l k} X _ {k i}\tag{25}
\]

in other words, the one-to-one correspondence (24) transforms products into products when all the products in question are defined (products of matrices of matrices being defined in the sense of no. 2, Remark); when the submatrices \(Z_{li}\) of the product YX are calculated thus, this product is said to be carried out ``in blocks''.

This name arises from the fact that, when \(I = [1, p]\) and \(K = [1, q]\) , the table representing the matrix X is envisaged as divided into ``blocks'' forming an ``array of matrices''

\[
\left( \begin{array}{c c c c} X _ {1 1} & X _ {1 2} & \ldots & X _ {1 p} \\ X _ {2 1} & X _ {2 2} & \ldots & X _ {2 p} \\ \cdot & \cdot & \cdot & \cdot \\ X _ {q 1} & X _ {q 2} & \ldots & X _ {q p} \end{array} \right)
\]

which is considered as a symbol denoting X when p and q are specific integers sufficiently small for this to be practicable.

\subsection{MATRIX OF A SEMI-LINEAR MAPPING}

Let A, B be two rings, \(\sigma: \mathbf{A} \to \mathbf{B}\) a homomorphism of A into B, E a right (resp. left) A-module with basis \((e_i)_{i \in I}\) and F a right (resp. left) B-module with basis \((f_k)_{k\in K}\) . Let \(u:\mathbf{E}\to \mathbf{F}\) be a semi-linear mapping relative to \(\sigma\) and \(u(e_i) = \sum_{k\in K}f_ku_{ki}\left(\text{resp.} u(e_i) = \sum_{k\in K}u_{ki}f_k\right)\) , where the \(u_{ki}\) are therefore elements of B; by definition, the matrix \(M(u) = (u_{ki})\) of type \(\mathbf{K}\times \mathbf{I}\) is also called the matrix of \(u\) with respect to the bases \((e_i)\) and \((f_k)\) . By the same calculation as in Proposition 2 of no. 4, it is immediately verified that for all \(x\in \mathbf{E}\) , if I and K are finite,

\[
M (u (x)) = M (u). \sigma (M (x))\tag{26 D}
\]

(resp.

\[
{ } ^ { t } M ( u ( x ) ) = \sigma ( { } ^ { t } M ( x ) ) . { } ^ { t } M ( u ) ) .\tag{26 G}
\]

Let C be a third ring, \(\tau: \mathbf{B} \to \mathbf{C}\) a homomorphism, G a right (resp. left) C-module with basis \((g_l)_{l \in \mathbf{L}}\) and \(v\) a semi-linear mapping of F into G relative to \(\tau\) ; if \(M(v)\) is the matrix of \(v\) with respect to \((f_k)\) and \((g_l)\) and \(M(v \circ u)\) the matrix of \(v \circ u\) relative to \((e_i)\) and \((g_l)\) , then, if I, K and L are finite,

\[
M (v \circ u) = M (v). \tau (M (u))\tag{27 D}
\]

(resp.

\[
{ } ^ { t } M ( v \circ u ) = \tau ( { } ^ { t } M ( u ) ) . { } ^ { t } M ( v ) ) .\tag{27 G}
\]

To show for example (27 D), note that for all \(x \in E\) , by (26 D),

\[
\begin{array}{r l} M (v \circ u) \cdot \tau (\sigma (M (x))) & = M (v (u (x))) \\ & = M (v) \cdot \tau (M (u (x))) = M (v) \cdot \tau (M (u)) \cdot \tau (\sigma (M (x))), \end{array}
\]

whence (27 D) by Proposition 1 of no. 3.

Suppose finally that \(\sigma: \mathbf{A} \to \mathbf{B}\) is an isomorphism; then recall that \(^t u: \mathbf{F}^* \to \mathbf{E}^*\) is a semi-linear mapping relative to \(\sigma^{-1}\) (§ 2, no. 5); when I and K are finite, the matrix \(^t u\) with respect to the dual bases \((f_k^*)\) and \((e_i^*)\) is given by

\[
M (^ {t} u) = \sigma^ {- 1} (^ {t} M (u))\tag{28}
\]

for, by definition, supposing for example that E and F are right modules, \(\langle^t u(f_k^*), e_i \rangle^\sigma = \langle f_k^*, u(e_i) \rangle\) when \(\sigma\) is denoted by \(x \mapsto x^\sigma\) .

Remark. Let A be a ring and \(\sigma\) an anti-endomorphism of A (no. 3); consider the two following situations:

(GD) E is a left A-module, F a right A-module and u a Z-linear mapping of E into F such that \(u(ax) = u(x)\sigma(a)\) for \(a \in A\) , \(x \in E\) ; in other words, u is a semi-linear mapping relative to \(\sigma\) of the right \(A^{0}\)-module E into the right A-module F.

(DG) E is a right A-module, F a left A-module and u a Z-linear mapping of E into F such that \(u(xa) = \sigma(a)u(x)\) for \(a \in A\) , \(x \in E\) ; in other words, u is a semi-linear mapping relative to \(\sigma\) of the left \(A^{0}\)-module E into the left A-module F.

In the two cases, the matrix \(M(u)\) of u relative to bases of E and F has its elements in A; if these bases are finite, then, for all \(x \in E\) , we have the respective formulae

\[
M (u (x)) = M (u). \sigma (M (x))\tag{17 GD}
\]

\[
{ } ^ { t } M ( u ( x ) ) = \sigma ( { } ^ { t } M ( x ) ) . { } ^ { t } M ( u ) ,\tag{17 DG}
\]

the products on the two sides being calculated in A. This follows immediately from (26 D) and (26 G) respectively.

\subsection{SQUARE MATRICES}

DEFINITION 4. A matrix whose rows and columns have the same indexing set is called a square matrix.

A square matrix with n rows and n columns is called a matrix of order n.

Remark. It should be noted that a matrix for which the indexing sets of the rows and columns have the same cardinal but are not identical, must not be considered as a square matrix; in particular, the product of two such matrices over a ring is not defined.

Clearly addition and multiplication of square matrices over A with a finite set as indexing set of the rows and columns, define on the set of these matrices a ring structure because of formulae (7), (8) and (9) (no. 3); the matrix \((\delta_{ij})\) , where \(\delta_{ij}\) is the Kronecker index (for \(i \in I\) , \(j \in I\) ), is the unit element of this ring and is denoted by \(I_n\) or \(1_n\) when I has \(n\) elements. When \(\mathbf{I} = [1, n]\) , the ring of matrices thus defined is denoted simply by \(\mathbf{M}_n(\mathbf{A})\) ; the group of invertible elements of \(\mathbf{M}_n(\mathbf{A})\) is denoted by \(\mathbf{GL}_n(\mathbf{A})\) or \(\mathbf{GL}(n, \mathbf{A})\) .

For a square matrix \(U = (a_{ij})\) of order n over A to be right (resp. left) invertible, it is necessary and sufficient that, for every system \((b_i)_{1 \leq i \leq n}\) of elements of A, the system of n equations in n unknowns

\[
\sum_ {j = 1} ^ {n} a _ {i j} x _ {j} = b _ {i} \quad (1 \leqslant i \leqslant n)
\]

\[
\left(\text {resp.} \sum_ {j = 1} ^ {n} x _ {j} a _ {j i} = b _ {i}\right)
\]

have one solution \((x_{i})\) in A.

Let I be a finite indexing set, A a ring and E a right (resp. left) A-module with basis \((e_{i})_{i\in I}\) . For every endomorphism u of E, the matrix \(M(u)\) of u with respect to the two bases identical with \((e_{i})\) is a square matrix; more briefly, it is called the matrix of u with respect to the basis \((e_{i})\) .

Suppose that \(\mathbf{I} = [1, n]\) . The mapping \(u \mapsto M(u)\) (resp. \(u \mapsto {}^t M(u)\) ) is an isomorphism of the ring \(\operatorname{End}_{\mathbf{A}}(\mathbf{E})\) onto \(\mathbf{M}_n(\mathbf{A})\) (resp. onto the opposite ring of \(\mathbf{M}_n(\mathbf{A})\) , as follows from formulae (18 D) (resp. 18 G)) (no. 4). The invertible elements of the ring \(\mathbf{M}_{n}(\mathbf{A})\) called invertible matrices, correspond under the mapping \(u \mapsto M(u)\) (resp. \(u \mapsto {}^{t}M(u)\) ) to the automorphisms of E; the group \(\mathbf{GL}(n, \mathbf{A})\) is therefore canonically identified with the group \(\mathbf{GL}(\mathbf{A}_{d}^{n})\) .

If \(u\) is an automorphism of \(\mathbf{E}\) , its contragredient \(\check{u}\) is an automorphism of the left (resp. right) A-module \(\mathbf{E}^*\) , such that \(\check{u} = (^{t}u)^{-1} = ^{t}(u^{-1})\) (§ 2, no. 5, Definition 6); if \(M(\check{u})\) is the matrix of \(\check{u}\) with respect to the dual basis \((e_{i}^{*})\) , then, by virtue of Proposition 3 (no. 4),

\[
M (\check {u}) = (^ {t} M (u)) ^ {- 1} = ^ {t} M (u ^ {- 1}).\tag{29}
\]

For every invertible matrix X, it therefore follows that \({}^{t}(X^{-1}) = ({}^{t}X)^{-1}\) ; this matrix is also denoted by \({}^{t}X^{-1}\) and called the contragredient of the matrix X.

Let \(\sigma\) be an automorphism of the ring A; for every semi-linear mapping \(u: \mathbf{E} \to \mathbf{E}\) relative to \(\sigma\) , the matrix \(M(u)\) of this mapping with respect to a basis \((e_i)\) of E is also a square matrix. It follows immediately from (27 D) (no. 6) that, if \(u\) is bijective, then

\[
M (u ^ {- 1}) = (\sigma^ {- 1} (M (u))) ^ {- 1}.
\]

Let E be an A-module which is the direct sum of a finite family \((\mathbf{E}_{i})_{i\in I}\) of submodules; for every endomorphism u of E, the matrix \(M(u)=(u_{kt})\) of u with respect to the two decompositions of E identical with \((\mathbf{E}_{i})\) (no. 5) is a square matrix of linear mappings. In order that \(u(\mathbf{E}_{i})\subset\mathbf{E}_{i}\) for all \(i\in I\) , it is necessary and sufficient that \(u_{kt}=0\) for \(k\neq i\) . When \(\mathbf{I} = [1, n]\) , the relations

\[
u (\mathrm{E} _ {i}) \subset \mathrm{E} _ {i} + \mathrm{E} _ {i + 1} + \dots + \mathrm{E} _ {n} \quad (1 \leqslant i \leqslant n)
\]

are equivalent to the relations \(u_{ki} = 0\) for \(k < i\) .

Examples of square matrices. I. Diagonal matrices. In a square matrix

\[
M = \left(m _ {\iota \kappa}\right) _ {(\iota, \kappa) \in I \times I},
\]

the elements both of whose indices are equal are called diagonal elements and the family \((m_{\iota})_{\iota \in I}\) is called the diagonal of \(M\) ; a square matrix \(M = (m_{\iota \kappa})\) over a ring, whose elements other than the diagonal elements are zero, is called a diagonal matrix. For every family \((a_{\iota})_{\iota \in I}\) of elements of a ring A, the diagonal matrix \((m_{\iota \kappa})\) such that \(m_{\iota} = a_{\iota}\) for all \(\iota \in I\) is denoted by \(\operatorname{diag}(a_{\iota})_{\iota \in I}\) (or \(\operatorname{diag}(a_1, a_2, \ldots, a_n)\) when \(I = [1, n]\) ). In the set \(\mathbf{M}_n(A)\) of square matrices of order \(n\) over A, the unit matrix \(I_n\) is a diagonal matrix and also every multiple \(aI_n = I_na\) of this matrix by a scalar \(a\) (the diagonal matrix (called scalar) all of whose diagonal elements are equal to \(a\) ).

For every family \((d_i)_{1\leqslant i\leqslant n}\) of elements of A and every matrix \(X = (x_{ij})\) of type \((n,q)\) (resp. \((p,n))\) over A, writing \(D = \mathrm{diag}(d_i)\) ,

\[
\left\{ \begin{array}{l} D X = (d _ {i} x _ {i j}) \\ X D = (x _ {i j} d _ {j}). \end{array} \right.\tag{30}
\]

In particular, for two diagonal matrices of order n,

\[
\begin{array}{c} \operatorname{diag} (a _ {i}) + \operatorname{diag} (b _ {i}) = \operatorname{diag} (a _ {i} + b _ {i}) \\ \operatorname{diag} (a _ {i}). \operatorname{diag} (b _ {i}) = \operatorname{diag} (a _ {i} b _ {i}). \end{array}\tag{31}
\]

The diagonal matrices therefore form a subring of \(\mathbf{M}_{n}(\mathbf{A})\) isomorphic to the product ring \(A^{n}\) ; the scalar matrices form a subring isomorphic to A.

\begin{enumerate}
\def\labelenumi{\Roman{enumi}.}
\setcounter{enumi}{1}
\tightlist
\item
  Permutation matrices; monomial matrices. Let \(\pi\) be any permutation of a finite set I and let \((e_i)_{i\in I}\) be the canonical basis of the A-module \(\mathbf{E} = \mathbf{A}_d^{\mathbf{I}}\) ; there exists one and only one endomorphism \(u_{\pi}\) of E such that, for all \(i\in I\) , \(u_{\pi}(e_i) = e_{\pi(i)}\) (§ 1, no. 11, Corollary 3 to Proposition 17). For all \(i\in I\) , the column of index \(i\) in the matrix \(M(u_{\pi})\) with respect to the basis \((e_i)\) has all its elements zero except the one in the row of index \(\pi(i)\) , which is equal to 1. By an abuse of language, \(M(u_{\pi})\) is called the matrix of the permutation \(\pi\) . It is immediate that for any two permutations \(\sigma, \tau\) of I, \(u_{\sigma \tau} = u_{\sigma} \circ u_{\tau}\) and that for the identity permutation \(\varepsilon, u_{\varepsilon}\) is the identity; the mapping \(\pi \mapsto M(u_{\pi})\) is therefore an isomorphism of the symmetric group \(\mathfrak{S}_{\mathbf{I}}\) onto the group of permutation matrices.
\end{enumerate}

Each row and each column of a permutation matrix contains only a single element \(\neq0\) . A finite square matrix R over a non-zero ring A, with this property, is called a monomial matrix; let \(r_{i}\) be the unique element \(\neq0\) in the column of R of index i and let \(\pi(i)\) be the index of the row where this element is; clearly \(\pi\) is a permutation of the indexing set I and \(R=M(u_{\pi})D\) , where \(D=\mathrm{diag}(r_{i})\) .

\begin{enumerate}
\def\labelenumi{\Roman{enumi}.}
\setcounter{enumi}{2}
\tightlist
\item
  Triangular matrices. In the ring \(\mathbf{M}_{n}(\mathbf{A})\) of square matrices of order n over a ring A, any matrix \((a_{ij})\) such that \(a_{ij}=0\) for i\textgreater j (resp. i\textless j) is called an upper (resp. lower) triangular matrix; it is also said that such a matrix has only zeros below (resp. above) its diagonal. It is immediately established that the upper (resp. lower) triangular matrices form a subring S (resp. T) of \(\mathbf{M}_{n}(\mathbf{A}),\mathbf{S}\cap\mathbf{T}\) being obviously the ring of diagonal matrices.
\end{enumerate}

The set S' (resp. T') of matrices in S (resp. T) whose diagonal elements are invertible is a multiplicative group of matrices called the upper (resp. lower) total triangular group, this follows immediately from § 1, no. 11, Remark 5. The set \(S_{1}\) (resp. \(T_{1}\)) of matrices in S (resp. T) whose diagonal elements are all equal to 1 is a subgroup of the above group, called the upper (resp. lower) strict triangular group, and every matrix M \ensuremath{\in} S' (resp. M \ensuremath{\in} T') whose diagonal is \((d_{i})\) , may be written as \(M = DM_{1} = M_{1}'D\) , where \(D = \mathrm{diag}(d_{i})\) and \(M_{1}\) and \(M_{1}'\) matrices belonging to \(S_{1}\) (resp. \(T_{1}\)).

\begin{enumerate}
\def\labelenumi{\Roman{enumi}.}
\setcounter{enumi}{3}
\tightlist
\item
  Diagonal and triangular matrices of matrices. Let \((\mathbf{I}_k)_{1 \leqslant k \leqslant p}\) be a partition of the finite set I; every square matrix over a ring A with indexing set I can be written in the form of a square matrix of matrices corresponding to the same partition \((I_{k})\) of the indexing set of the rows and the indexing set of the columns (no. 5)
\end{enumerate}

\[
\left( \begin{array}{c c c c} X _ {1 1} & X _ {1 2} & \ldots & X _ {1 p} \\ X _ {2 1} & X _ {2 2} & \ldots & X _ {2 p} \\ \cdot & \cdot & \cdot & \cdot \\ X _ {p 1} & X _ {p 2} & \ldots & X _ {p p} \end{array} \right)\tag{32}
\]

where each \(X_{kk}\) is a square matrix with \(I_{k}\) as indexing set of the rows and columns.

With this notation, (32) will be called a diagonal (resp. upper triangular, resp. lower triangular) matrix of matrices if all the matrices \(X_{ij}\) such that \(i \neq j\) (resp. i \textgreater{} j, resp. i \textless{} j) are zero. The interpretation of endomorphisms u whose matrix is a diagonal, resp. triangular, matrix of matrices has been seen earlier, by considering the corresponding matrix \(M(u)\) of linear mappings. The lower triangular (resp. upper triangular, diagonal) matrices of matrices for a given partition \((\mathbf{I}_{k})\) of I form subrings of the ring of matrices \(A^{I \times I}\) . In particular, the ring of diagonal matrices of matrices relative to

the partition \((\mathbf{I}_k)\) is isomorphic to the product \(\prod_{k=1}^{p} \operatorname{End}_{\mathbf{A}}(\mathbf{E}_k)\) .
